\section{A.1. What the section 6 maps can and cannot be}

We use labelled markings and work over $\C$. Unless explicitly stated otherwise,
$2g-2+n>0$. A stable graph $\Gamma$ specifies the connected components of the
normalization, the paired half-edges at nodes, and the external markings. On
67/075 the propagator at a node with branches $p_+,p_-$ is explicitly a nonzero
vector in $T_{p_+}Y\otimes_\C T_{p_-}Y$, not an infinitesimal equivalence class
of a specified analytic plumbing family. The unrestricted target therefore
contains every nonzero vector in each of these lines. This distinction is
decisive.

\begin{definition}[Conformal shadow]\label{A:shadowdef}
\status{Proved well-defined; convention checked against 67/062--064.}
For a holomorphically embedded pointed disc $(D,p)$ in a Riemann surface $Y$,
choose a biholomorphism $f:\D\to D^\circ$ with $f(0)=p$. Its shadow is
$df_0(\D)\subset T_pY$. In a local coordinate $z$ with $z(p)=0$, its radius
is $r_z(D,p)=|(z\circ f)'(0)|$. Changing $f$ by a rotation changes neither
the shadow nor its radius. Changing coordinate to $w$ multiplies the radius
by $|dw/dz(p)|$.
\end{definition}

\begin{lemma}[The tensor produced by a toric gluing]\label{A:tensor}
\status{Proved.}
Let $f_\pm:(\D,0)\to(D_\pm,p_\pm)$ be pointed conformal parametrizations.
An orientation-reversing isometry between their boundary circle torsors has
the form $\zeta_- = e^{i\theta}/\zeta_+$ on $|\zeta_+|=1$. The induced
propagator is
\[
 q=e^{i\theta}df_{+,0}(1)\otimes df_{-,0}(1)
       \in T_{p_+}Y\otimes T_{p_-}Y,
 \qquad |q|_{z_+,z_-}=r_{z_+}(D_+,p_+)r_{z_-}(D_-,p_-).
\]
The inverse tensor, viewed as a nonzero bilinear pairing into $\C$, has
reciprocal modulus.
\end{lemma}
\begin{proof}
The associated Hermitian norms identify both boundary torsors with unit
circles. Every equivariant orientation-reversing identification between them
has the displayed form. On replacing $f_\pm(\zeta)$ by
$f_\pm(e^{i\alpha_\pm}\zeta)$, the parameter changes to
$e^{i(\theta-\alpha_+-\alpha_-)}$, while the two derivatives acquire the
factors $e^{i\alpha_+}$ and $e^{i\alpha_-}$. Their product is unchanged.
This proves both the intrinsic formula and the modulus assertion. It also
checks the dual convention mentioned in the margin of 67/078.
\end{proof}

\begin{lemma}[A radius bound on a three-pointed sphere]\label{A:koebe}
\status{Known theorem, applied with all normalizations specified.}
If $f:\D\hookrightarrow\PP^1\setminus\{1,\infty\}$ is holomorphic,
$f(0)=0$, and $r=|f'(0)|$, then $r\leq4$.
\end{lemma}
\begin{proof}
The normalized schlicht function $F=f/f'(0)$ omits $1/f'(0)$. Koebe's
quarter theorem \citep[Appendix, pp.~10--13]{C_Siu2005} gives $\{|w|<1/4\}\subset F(\D)$, so
$|1/f'(0)|\geq1/4$. Thus $r\leq4$. Only this non-strict inequality is
needed below; no boundary regularity or equality-case assertion is used.
\end{proof}

\begin{theorem}[The key case $(0,4)$ refutes the global isomorphism]\label{A:04}
\status{Refuted, with a complete geometric counterexample.}
For the targets described on 67/075--078, the capping-and-forgetting maps
$MD_{0,4}\to M_{\mathrm{st}}N_{0,4}$ and, whenever defined with the same
unrestricted nodal tensor, $MDB_{0,4}\to M_{\mathrm{st}}ND_{0,4}$ are not
surjective. In particular, they are not algebraic, complex analytic,
real analytic, topological, or set-theoretic isomorphisms onto those targets.
Lead L2 is false.
\end{theorem}
\begin{proof}
Take two copies $Y_+,Y_-$ of $\PP^1$, with coordinates $z_+,z_-$. Identify
$0\in Y_+$ with $0\in Y_-$. Give each component its two distinct labelled
external marks $1$ and $\infty$. The resulting connected curve $Y$ has
arithmetic genus zero and four smooth distinct markings; each rational
component has three special points, so it is stable. Its labelled
normalization has no nonidentity automorphism. Decorate its unique node with
\[
 q_0=17\,\frac{\partial}{\partial z_+}\otimes
             \frac{\partial}{\partial z_-}.
\]
This is a legitimate nonzero tangent tensor, and hence a target object.

Suppose it had a source object. Cutting at the chosen separating circle
and capping would give pointed discs $(D_+,0)$ and $(D_-,0)$ in these
normalized spheres. Each disc avoids both external marks. This remains
true if the external boundary components themselves have positive size:
their cap centres are outside the node cap. Lemma~\ref{A:koebe} therefore
gives $r_+\leq4$ and $r_-\leq4$. Lemma~\ref{A:tensor} gives
$|q_0|=r_+r_-\leq16$, a contradiction to $|q_0|=17$.

For the bordered target choose, on each component with the node regarded
as a cusp, any two positive-length admissible external caps. Such caps
exist by the pants construction, equivalently by
Theorem~\ref{A:grafting}. The tensor $q_0$ remains independent additional
data and the same contradiction applies. If a different definition of
``admissible external caps'' imposes compatibility with the omitted node
caps, that is a restriction of the target and must be stated explicitly;
it is not the unrestricted target above. In the reciprocal pairing
convention, replace $17$ by $1/17$. There is again no preimage.
\end{proof}

\begin{theorem}[Obstruction on every nodal stratum]\label{A:alltypes}
\status{Proved.}
Fix any stable pointed complex curve $Y$ with at least one node. The
propagators obtainable from embedded caps on its fixed normalization are
bounded in each nodal tangent line, for any fixed Hermitian norms. Thus
the unrestricted nonzero-tensor target over $Y$ is never the image of
capping. For every stable type with $3g-3+n>0$ this gives a counterexample
to a global isomorphism. The sole stable type of complexity zero is
$(0,3)$; there are no nodal stable curves of that type, and capping with
all three external discs retained is an equivalence.
\end{theorem}
\begin{proof}
We bound $df_0$ for any embedded disc $f:(\D,0)\hookrightarrow(Z,p)$ on
a fixed compact normalization component $Z$, avoiding its other special
points.

If $Z$ has genus zero, stability gives at least two other special points.
Put them at $1,\infty$ and $p$ at $0$. Lemma~\ref{A:koebe} gives a finite
bound, which transfers to any other fixed tangent norm by the derivative
of this coordinate change.

If $Z=\C/\Lambda$ has genus one, choose a flat metric of area $A$. Lift
$f$ to $F:\D\to\C$. Because $f$ is an embedding, change of variables
gives $\int_\D|F'|^2\,dx\,dy\leq A$. Orthogonality of the powers of
$z$ in the Taylor expansion $F(z)=\sum a_jz^j$ gives
$\pi|a_1|^2\leq\sum_{j\geq1}\pi j|a_j|^2\leq A$.
Thus $|df_0|\leq\sqrt{A/\pi}$ in that metric.

If $Z$ has genus at least two, equip it with its curvature $-1$ metric.
Schwarz--Pick applied to the lifted map into the disc gives
$|df_0|_Z\leq2$, since the Poincare density of $\D$ at zero is $2$.

Multiplying the bounds on the two branches proves the assertion for
each node. A finite automorphism group cannot remove the obstruction:
average a Hermitian norm over that group, or take the maximum of the
finitely many transformed bounds. Choose a tensor exceeding that bound.

For $g=0,n\geq4$, split the external marks into sets of sizes $2$ and
$n-2$ and join the corresponding rational components at one node. For
$g\geq1$ and a stable type other than the excluded $(1,0)$, identify two
additional distinct points on a smooth genus-$g-1$ curve with $n$ marked
points. Its normalization is stable, since
$2(g-1)-2+(n+2)=2g-2+n>0$. These constructions produce the required
nodal curves. In particular $(1,1)$ has a rational normalization with
the two node branches and one marking; the genus-zero argument applies.
For $(0,3)$, the dimension/stability count forces no nodes. Removing
the retained three admissible discs and recapping are inverse operations.
\end{proof}

\begin{proposition}[Why arbitrary smooth cuts require more data]\label{A:arbitrarycuts}
\status{Proved under stated interpretation; the hypothesis matters.}
If the source on 67/074 includes arbitrary smooth cut circles, the
canonical hyperbolic boundary torsors need not be respected by the
identification of their two banks. In that interpretation the asserted
pointed-disc anti-isomorphism, and hence the asserted tensor functor,
is not defined by the stated construction. Restricting to geodesic
cuts in the canonical hyperbolic metric removes this problem but does
not remove Theorem~\ref{A:04}.
\end{proposition}
\begin{proof}
Take two hyperbolic pants with one boundary of the same positive
length and use arclength parameters $t$ on the corresponding circles.
Glue them by the orientation-reversing real analytic map
$t\mapsto-t-\varepsilon\sin t$, where $0<\varepsilon<1$.
Its derivative never vanishes, so analytic conformal welding produces
a surface with an analytic cut circle. On cutting it again, we recover
the original conformal pants and hence their original canonical
hyperbolic metrics. A pointed anti-holomorphic isomorphism between
their cap discs restricts to $t\mapsto c-t$ on their normalized
boundaries. The chosen gluing is not of that form. Thus it does not
produce the anti-isomorphism asserted on 67/074--075.

For a geodesic cut, restriction of the original metric to each piece
has curvature $-1$ and geodesic boundary. Uniqueness of that metric
identifies it with the piece's canonical metric. The two bank lengths
and arclength torsors then agree, so Lemma~\ref{A:tensor} applies.
\end{proof}

\begin{theorem}[The exact surgery equivalence retaining the caps]\label{A:fullcaps}
\status{Proved.}
For a fixed stable graph $\Gamma$, consider the following two groupoids:
(i) hyperbolic bordered surfaces with labelled disjoint essential
geodesic cuts of graph $\Gamma$, with conformal isomorphisms preserving
all labels; (ii) their compact normalized components with a full
admissible system of pointed cap discs at every half-edge and external
mark, equal hyperbolic lengths on paired half-edges, and a specified
orientation-reversing toric identification on each paired boundary.
Cutting and capping gives an equivalence of these groupoids, with
inverse obtained by deleting the node caps and welding the paired
boundaries. Cuspidal external components are allowed as zero discs.
\end{theorem}
\begin{proof}
Admissibility means exactly that the circle torsor obtained from the
pointed cap agrees with arclength for the canonical hyperbolic metric
on the complement of all cap interiors. Equal paired lengths make a
toric identification an isometry. In Fermi coordinates near a
geodesic boundary the metric is
$ds^2+\cosh^2(s)\,dt^2$, with $t$ periodic of the specified length.
Identification of $s=0$ with the reverse boundary orientation gives
this same smooth metric on both sides, since $\cosh$ is even. It has
curvature $-1$ and its seam is geodesic. Its conformal structure is
therefore the unique analytic welding structure. The metric remains
complete at any cusps.

Cutting the welded surface recovers the pieces and their isometric
boundary identifications. Conversely, cutting and capping a source
surface and then performing this welding recovers its original
metric and conformal structure. The normalized disc parametrizations
are unique up to rotations; the torsor formulation removes exactly
that ambiguity. An isomorphism on either side extends or restricts
uniquely on all pieces, which proves full faithfulness as well as
essential surjectivity. Forgetting the node discs in this equivalence
is the extra operation whose unrestricted inverse fails above.
\end{proof}

\subsection{The precise target on a fixed cut stratum}

The next result gives the radius restriction for every stable type and
computes the two maps proposed on 67/081--082. We use the radius theorems
proved in Section~\ref{C:A2b-energy-section} and the subsequent image
theorem. Thus the restriction of the target is determined intrinsically,
not by an unspecified ``sufficiently small'' condition. Here a labelled
cut surface means the moduli-groupoid object with its cut circles and
punctures labelled, not a surface with an additional global topological
marking. In the fully marked Teichmuller space a boundary twist is real,
not circle-valued. One then lifts each tensor phase to its universal
cover and obtains the corresponding statements there. Thus the
integral Dehn twists are quotiented on both sides, or retained on both
sides; they are not silently lost from only one side.

\begin{theorem}[Proved: the exact propagator domain and the tangent map]
\label{A:graph-radius-domain}
Let $\Gamma$ be a connected stable graph of type $(g,n)$, with edge set
$E\ne\varnothing$. Fix its pointed compact normalization components
$X_v$, with every branch and external marking labelled. For an ordinary
fibre, retain an identification of the normalization with these fixed
objects; merely labelling the special points need not remove their
finite automorphisms. All constructions below are equivariant for
those automorphisms and also define the corresponding groupoid maps.
We first take
all external ends to be cusps. At each special point choose a local
coordinate, used only to measure radii and tensor moduli. Let $I_v$ be
the set of node branches on $X_v$. For $a\in[0,\infty)^E$, define
$(B_va)_h=a_e$ when the branch $h$ belongs to the edge $e$; both branches
of a loop are included separately. Let $H_v$ be the Jenkins--Strebel
support function for these active points, with all other special points
kept marked, and put
\[
 H_\Gamma(a)=\sum_v H_v(B_va),\qquad H_v(0)=0.
\]
Define the explicit open convex set
\[
 K_\Gamma=\left\{y\in\R^E:
       \sum_e a_e y_e<H_\Gamma(a)
       \text{ for every }a\geq0,\ \sum_ea_e=1\right\}.
\]
The attainable nonzero node tensors are exactly those with
$(\log|q_e|)_e\in K_\Gamma$. Every such tensor system has a unique
labelled geodesic-cut preimage, up to the equivalence of
Theorem~\ref{A:fullcaps}. More precisely, if
$x_e=(\ell_e/2\pi)^2$, the logarithmic modulus map is
\[
 G_\Gamma(x)=\sum_v B_v^T F_v(B_vx),\qquad
 G_\Gamma:(0,\infty)^E\xrightarrow{\ \sim\ }K_\Gamma,
\]
where $F_v$ is the vector of logarithmic cap radii on $X_v$.
It is a real analytic diffeomorphism with symmetric positive definite
derivative. In the phase coordinates specified by normalized cap maps,
the complete tangent map is block diagonal with blocks $DG_\Gamma$
and the identity on $\R^E$. This is the tangent map on the fixed
normalization fibre. When the normalization components also vary,
the differential is block triangular with the identity on the
normalization-moduli base and this invertible fibre block. In
particular the full real tangent map is invertible everywhere on
the fixed graph stratum.
\end{theorem}

\begin{proof}
For positive $x$, Mondello's inverse provides on each $X_v$ the unique
cap system with branch lengths $2\pi\sqrt{x_e}$. The two lengths paired
at an edge agree by construction. Lemma~\ref{A:tensor} gives exactly the
displayed formula for $G_\Gamma$, since the logarithm of a product of
the two radii is their sum. Conversely every geodesic-cut source has
these equal branch lengths, and hence is obtained by this construction.

The chain rule and the energy identity give
\[
 DG_\Gamma(x)=\sum_v B_v^T DF_v(B_vx)B_v.
\]
For any nonzero $u\in\R^E$, some branch of an edge with $u_e\ne0$
has nonzero coordinate in $B_vu$. The corresponding positive definite
summand makes $u^TDG_\Gamma(x)u>0$. Vertices without node branches
contribute zero and cause no exception. Integrating this inequality
along line segments proves injectivity, and the inverse function
theorem gives a local real analytic inverse.

The componentwise strict support inequalities imply
$G_\Gamma(x)\in K_\Gamma$: a nonzero $a$ has at least one nonzero
$B_va$, at which the inequality is strict. The finite-area support
gap on the components also gives the uniform bound
\[
 0<H_\Gamma(x)-x\cdot G_\Gamma(x)
       \leq\sum_v(2g_v-2+n_v)=2g-2+n.
\]
Here $n_v$ counts all node branches and external markings. The last
identity follows from $g=\sum_vg_v+|E|-|V|+1$ and
$\sum_v n_v=2|E|+n$.

We give the properness argument explicitly. Suppose $G_\Gamma(x_m)$
converges to $y\in K_\Gamma$. If $s_m=\sum_e x_{m,e}$ tends to
infinity along a subsequence, pass to another subsequence with
$a_m=x_m/s_m\to a$, $a\geq0$, $\sum a_e=1$. Homogeneity and the gap
give $H_\Gamma(a_m)-a_m\cdot G_\Gamma(x_m)\to0$.
Continuity of the support functions therefore gives
$H_\Gamma(a)=a\cdot y$, contradicting $y\in K_\Gamma$.
If a bounded subsequence has $x_{m,e}\to0$, the uniform collar bound
at either branch makes its log radius tend to $-\infty$. The other
branch radius has a fixed upper bound on its fixed normalization
component, so $(G_\Gamma(x_m))_e\to-\infty$, again a contradiction.
The inverse image of every compact subset of $K_\Gamma$ is therefore
compact in the positive orthant. Its image is nonempty, open, and
closed in the connected set $K_\Gamma$, and hence equals $K_\Gamma$.
This proves the asserted global inverse.

Choose the cap parametrizations so that their derivatives in the
specified local coordinates are positive real numbers. Each edge's
gluing phase $e^{i\theta_e}$ is then precisely the phase of $q_e$.
The source coordinates $(x,\theta)$ therefore map to
$(G_\Gamma(x),\theta)$. Their differential has the claimed blocks.
The welding equivalence supplies the surface and its uniqueness.
Changing coordinate norms translates the logarithmic target and its
support inequalities by the same vector, and changing coordinate
arguments rotates the phase torsors. Thus the intrinsic statement
does not depend on these auxiliary choices.
\end{proof}

\begin{corollary}[Proved: the fundamental-group test does not establish surjectivity]
\label{A:graph-pi-one}
For a fixed labelled normalization with its identification retained,
the map on ordinary fundamental groups
from the geodesic-cut source to the unrestricted tensor target is
the identity $\Z^E\to\Z^E$ in the twist generators, although its
image is the proper domain of Theorem~\ref{A:graph-radius-domain}.
The analogous statement holds for orbifold fundamental groups when
the normalization components vary and their finite automorphisms
are retained: inclusion of the restricted tensor domain is a
homotopy equivalence on each fixed graph stratum. This is compatible
with the everywhere invertible real tangent map. Neither test supplies
the missing properness onto the unrestricted target.
\end{corollary}

\begin{proof}
At a fixed normalization with its identification retained, take
$y_0=G_\Gamma(\mathbf 1)$. Linear
contraction in logarithmic radii to $y_0$ is a deformation retraction
of both $K_\Gamma\times(S^1)^E$ and
$\R^E\times(S^1)^E$ onto $\{y_0\}\times(S^1)^E$; convexity keeps
the former contraction inside $K_\Gamma$. Inclusion is therefore a
homotopy equivalence and preserves every circle generator. A full
turn in the normalized boundary twist is exactly a full turn in the
tensor phase, proving the assertion about the displayed generators.

Over varying normalization components use, intrinsically, the modulus
of the tensor produced when all node lengths equal $2\pi$ to define
the radial centre. This is a positive norm in each nodal tensor line;
it does not choose a phase in that line. The radial homotopies just
described are consequently defined without a local coordinate or
trivialization. Real analytic dependence in the component moduli
makes them continuous. They commute with relabellings and finite
automorphisms, and hence descend to the moduli orbifold, or equally
to its groupoid classifying space. The same deformation retraction
on both sides proves the orbifold assertion.

Properness fails for the larger target: take $x_m=t_m\mathbf1$ with
$t_m\to\infty$. The Jenkins--Strebel limit on every component shows
that $G_\Gamma(x_m)$ converges to a finite boundary point of
$K_\Gamma$, with equality in its positive diagonal support
inequality. Keeping all phases fixed gives a convergent sequence
of target tensors with no convergent source sequence. Thus the
inference from an invertible tangent map and a fundamental-group
isomorphism to a global isomorphism has a concrete failed hypothesis.
\end{proof}

\begin{example}[Proved: the entire key fibre in type $(0,4)$]
\label{A:04-exact-image}
Take the two labelled three-pointed spheres used in
Theorem~\ref{A:04}, with all four external ends cuspidal. Each node
cap has the same radius $r(\lambda)$, which increases strictly from
$0$ to $4$ as $\lambda$ increases from $0$ to infinity. The precise
tensor image in the specified coordinates is
\[
                    0<|q|<16.
\]
The parameter relation is $|q|=r(\lambda)^2$ and its phase is the
boundary twist. Thus this fibre is a punctured disc, whereas the
proposed unrestricted target fibre is $\C^*$. The derivative never
vanishes and both fundamental groups are $\Z$, with degree one map.
\end{example}

\begin{proof}
For one active puncture at zero on the three-pointed sphere,
the extremal Jenkins--Strebel domain is
$\C\setminus[1,\infty)$ and has conformal radius $4$, as the explicit
map $z=4u/(1+u)^2$ shows. The one-dimensional image theorem therefore
gives $F((0,\infty))=(-\infty,\log4)$, and the energy identity gives
$F'>0$. The two components are identically normalized, so their
radii agree at the common boundary length. Multiplication gives
exactly the claimed interval of tensor moduli. The assertions about
the derivative and fundamental group are also immediate from the
two preceding results.
\end{proof}

\begin{proposition}[Proved: the bordered image with all compatibility data retained]
\label{A:bordered-radius-algorithm}
The same finite-dimensional inversion determines the image when
external caps are retained. Add one positive variable $x_i$ for each
positive external boundary, and one output coordinate $\log r_i$
for its shadow. In $B_v$, copy a node variable to each of its two
flags and an external variable to its single flag. Then
$B^T(F_v)_v$ is a real analytic diffeomorphism onto the strict
support domain with support $\sum_vH_v(B_va)$.

For specified normalization components, node tensors, and actual
external pointed discs, the following is a necessary and sufficient
test for a geodesic-cut preimage: their tensor moduli and external
shadow radii belong to that domain; the unique recovered length
vector produces exactly the specified external discs. When both
conditions hold, the preimage is unique. Zero external radii are
handled on their corresponding face. This test states the additional
disc compatibility instead of assuming that forgetting node caps
preserves admissibility.
\end{proposition}

\begin{proof}
The matrix $B$ is still injective. The Hessian, support, finite-area
gap, and collar arguments in Theorem~\ref{A:graph-radius-domain}
therefore apply word for word, with a single cap radius instead of a
product at each external flag. This proves the indicated
diffeomorphism. A source necessarily gives its unique inverted
length vector and its canonical caps, proving necessity of both
tests. Conversely that length vector gives equal node-boundary
lengths, and the phase of each specified tensor gives its unique
toric gluing. The second test identifies the remaining external
caps with those required. Theorem~\ref{A:fullcaps} then supplies
the preimage. Its uniqueness follows from the radius inversion
and the uniqueness of toric gluing with its given phase.
\end{proof}


\subsection{Local plumbing and the limits of the proposed identification}

\begin{proposition}[The genuine local statement]\label{A:plumbing}
\status{Known, with the comparison explicitly limited to germs.}
For a stable complex curve $Y$ and a multicurve $\Gamma$ representing
its nodes, choose holomorphic coordinates $z_e,w_e$ near the two
branches of each node and local parameters $u$ for the pointed
normalization. For sufficiently small $t_e$, the equations
$z_ew_e=t_e$ yield local analytic deformation coordinates after the
appropriate Dehn-twist quotient. Their normal tangent directions are
canonically the lines $T_{p_e^+}Y\otimes T_{p_e^-}Y$.
This does not identify an entire product of punctured tangent lines
with the global cut-surface moduli.
\end{proposition}
\begin{proof}
The local coordinate assertion is \citet[Theorem 9.11 and Theorem
10.1]{HubbardKoch2014}. Their Section 8 specifies the plumbing family;
the theorem is invoked only on the neighbourhood where it applies.
For the tangent-line identification, replace the coordinates by
$z'_e=a_ez_e+O(z_e^2)$ and $w'_e=b_ew_e+O(w_e^2)$.
To first order normal to $t_e=0$, $t'_e=a_eb_et_e$, while
$\partial_{z'_e}\otimes\partial_{w'_e}
=(a_eb_e)^{-1}\partial_{z_e}\otimes\partial_{w_e}$.
Thus $t_e\partial_{z_e}\otimes\partial_{w_e}$ is intrinsic to first
order. Higher coordinate jets affect finite plumbing, so this
calculation supplies a normal-line identification, not a canonical
global inverse. Hubbard--Koch explicitly deny global injectivity of
their plumbing map immediately after Theorem 9.11. Independently,
Theorem~\ref{A:04} rules out the global inverse requested in L2.
\end{proof}

\begin{proposition}[Proved: the plumbing equation for the corrected inverse]
\label{A:cap-plumbing-equation}
The inverse on the restricted propagator domain in
Theorem~\ref{A:graph-radius-domain} is a plumbing in the canonical cap
coordinates, which themselves depend on the recovered length vector.
If $z_\pm$ are coordinates at the two node branches and
$f_\pm:\D\to D_\pm$ are normalized so that
$(z_\pm\circ f_\pm)'(0)=r_\pm>0$, put
$Z_\pm=r_\pm f_\pm^{-1}$. Its seam identification is exactly
\[
                       Z_+Z_-=q,\qquad
             q=r_+r_-e^{i\theta}.
\]
This identity does not identify the radius inverse with plumbing in
one fixed pair of coordinate germs for all $q$.
\end{proposition}

\begin{proof}
The canonical cap boundary is analytic, so its normalized conformal
coordinate extends univalently to a collar across the boundary.
On that boundary the toric identification is
$f_-^{-1}=e^{i\theta}/f_+^{-1}$, proving the displayed equation.
Take slightly enlarged coordinate discs of radii $R_\pm>r_\pm$.
They have $R_+R_->|q|$, so the transition $Z_-=q/Z_+$ identifies
overlapping annuli. Its restriction to the seam is the specified
identification; its holomorphic extension gives the same conformal
welding by uniqueness. The derivatives of $Z_\pm$ relative to
$z_\pm$ at the origins are one, so the coefficient is the actual
tensor coefficient used in Lemma~\ref{A:tensor}.

The inverse radius theorem first recovers the lengths, and only then
these cap maps. Their dependence is real analytic on each face.
Higher jets of $Z_\pm$ therefore move with $q$. The tensor supplies
the first derivatives and product, whereas a fixed-coordinate
plumbing also specifies those higher jets. Equality of the equations
in their respective moving coordinates is consequently the precise
comparison that the construction proves.
\end{proof}

\begin{proposition}[Known: the applicable horocyclic and metric comparisons]
\label{A:plumbing-sources}
Kra's horocyclic construction and Wolpert's hyperbolic plumbing
are valid comparisons after their coordinate choices and domain
restrictions are retained. They do not parametrize arbitrary
nonzero normal tensors by the manuscript's canonical capping map.
\end{proposition}

\begin{proof}
For a maximal pants graph with $d=3g-3+n$, \citet[Theorems~1
and~3, pp.~501--503; \S\S2--3]{Kra1990} uses cusp coordinates and
the equation $zw=t$. His tame construction contains the polydisc
$|t|_{\max}<e^{-2\pi}$. The full domain $D_\Gamma$ is a domain of
holomorphy satisfying
\[
  (\D_{e^{-2\pi}})^d\subset D_\Gamma
                    \subset(\D_{e^{-\pi/2}})^d,
\]
and its map to ordinary moduli can have infinitely many preimages.
The logarithmic coordinates on the fully marked Teichmuller space
have Dehn twists $\tau_j\mapsto\tau_j+2$.
Thus his domain and marking quotient must both be kept.
The present inverse instead chooses its discs by radius inversion.
Proposition~\ref{A:cap-plumbing-equation} proves the shared local
equation without identifying these different coordinate choices.

For the hyperbolic comparison, use normalized cusp coordinates
$z=e^{2\pi i\zeta}$ and $w$ on the limiting components. On
$|t|<|z|<1$, the complete annular metric is
\[
 ds_t^2=\left(
 \frac{\pi}{\log|t|}
 \csc\frac{\pi\log|z|}{\log|t|}
 \left|\frac{dz}{z}\right|\right)^2.
\]
Integrating on $|z|=|t|^{1/2}$ gives core length
$2\pi^2/(-\log|t|)$. Wolpert's curvature correction transfers this
model to the plumbed surface, with normalized parameter
$\tau_j$ including the first derivatives of the chosen coordinates:
\citet[Expansion~4.2, p.~445, and Example~4.3, p.~446]{Wolpert1990}.
In particular, for one pinched curve and fixed remaining data,
\[
 \ell(t)=\frac{2\pi^2}{-\log|t|}
                   +O((-\log|t|)^{-2}).
\]
The coordinate scaling changes the logarithm by a bounded constant,
which explains this form of the error.
The punctured-surface version and its hypotheses are stated in
\citet[\S2, Theorems~4 and~6]{ObitsuWolpert2008}: the local
parameters are real analytic in the normalization variables and
holomorphic in $t$ with those variables fixed. Here their metric
interpolation called grafting is the construction in that paper,
distinct from attaching the infinite flat cap cylinders.

Finally Hubbard--Koch's local chart theorem, applied in
Proposition~\ref{A:plumbing}, identifies these plumbing germs with
the analytic Deligne--Mumford deformation problem. The first-order
transition $t'=ab\,t$ proves the common normal-line interpretation.
The independent unrestricted-vector counterexample in
Theorem~\ref{A:04} excludes any global inverse of the kind proposed
in L2, regardless of which of these local coordinate theorems is
used.
\end{proof}

\begin{proposition}[Known: the proved local Earle--Marden comparison]
\label{A:EarleMarden-local}
On a neighbourhood of a fixed nodal surface, the Earle--Marden
Kleinian plumbing family and the Hubbard--Koch family give the same
analytic marked deformation germ. This comparison follows from the
proved Appendix of \citet[pp.~41--42 in the opened author version]
{HubbardKoch2014}. It is a local comparison with a specified
Dehn-twist quotient.
\end{proposition}

\begin{proof}
We spell out the identification used here. Normalize each peripheral
parabolic of the nodal Kleinian group to $\zeta\mapsto\zeta+1$.
The added translation by $\tau_e$, with imaginary part exceeding a
fixed collar bound, opens that node. Passing to
$t_e=e^{2\pi i\tau_e}$ removes exactly the integral Dehn-twist
ambiguity. With the transition $\zeta_+=\zeta_-+\tau_e$, set
$z_e=e^{2\pi i\zeta_+}$ and $w_e=e^{-2\pi i\zeta_-}$.
Then $z_ew_e=e^{2\pi i\tau_e}=t_e$ exactly; in particular the
smoothing exponent is one, after these fixed normalizations.
Together with the parameters $u$ of the normalization components,
this constructs the $\Gamma$-marked analytic stable family on
$U\times\D_\rho^E$ described in the cited Appendix.

By Hubbard--Koch's Theorem~10.1 its universal property supplies
an analytic classifying map to their $\mathcal Q_\Gamma$.
At the nodal point, the $u$-directions are the same normalization
deformations on both sides; each $t_e$ opens just its own node
and has its nonzero normal tangent coefficient in
$T_{p_e^+}\otimes T_{p_e^-}$. Thus the derivative is an
isomorphism, by the independent smoothing directions and
normal-line calculation in Proposition~\ref{A:plumbing}.
The analytic inverse function theorem identifies the germs, and
the classifying property identifies their families. This is the
comparison proved by the Appendix, not a claim that the
neighbourhood can be chosen to be the whole tensor bundle.

The bibliography separately marks the original 2012
Earle--Marden chapter as not opened. We have verified the
mathematical comparison in the opened primary proof just
specified. The prompt does not identify the 2012 chapter as the
required source: the opened Appendix is itself a primary proof of the
requested construction comparison. The unavailable chapter remains a
bibliographic coverage limitation, and no statement is accepted on its
unopened authority.
\end{proof}

\begin{example}[Refuted: arbitrary fixed plumbing is everywhere a local chart]
\label{A:hinich-explicit}
Even a holomorphic plumbing of two three-pointed spheres with
disjoint coordinate discs can have zero derivative at a nonzero
parameter. We give exact rational data for the mechanism of
\citet[\S\S2.1--2.3]{Hinich2010}, rather than leaving its
parameters existential.
\end{example}

\begin{proof}
Set
\[
 R=\frac32,\quad B=\frac34,\quad C=10,\quad
 A=\frac{420}{569},\quad\varepsilon=\frac1{100},\quad
 \phi(\zeta)=(\zeta-R)^2.
\]
The map $\phi$ is injective on $|\zeta|<1+\varepsilon$:
two different points with the same image would sum to $2R=3$,
whereas their sum has modulus below $2(1+\varepsilon)<3$.
Its derivative has no zero there. Put
$a=\phi(A)$, $b=\phi(B)$, $c=\phi(C)$. The distinct marks
$a,b$ lie on the bounded side of
$\phi(\{1-\varepsilon<|\zeta|<1+\varepsilon\})$.
The marks $c,\infty$ lie on the other side, since both roots
$C,2R-C$ of $\phi(\zeta)=c$ have modulus above $1+\varepsilon$.

We specify the two caps to avoid any orientation ambiguity.
Attach a disc with centre $\zeta=\infty$, using the region
$|\zeta|>1-\varepsilon$, to the bounded-side component and its
annular collar. This makes the first sphere, marked by $a,b$
and the new centre. Attach the disc $|\zeta|<1+\varepsilon$,
with centre zero, to the unbounded-side component and its collar.
This makes the second sphere, marked by $c,\infty$ and its
new centre. The overlaps are identified by $\phi$.
Each compactified surface is a topological sphere with a
holomorphic atlas, hence is a Riemann sphere. The respective
node coordinates are
\[
             x=(1-\varepsilon)/\zeta,\qquad
             y=\zeta/(1+\varepsilon).
\]
Their unit discs contain none of the other markings, and lie
on distinct components, so the disjointness condition holds.
Plumbing $xy=t$ at
$t_0=(1-\varepsilon)/(1+\varepsilon)=99/101$
recovers the sphere with marks $a,b,c,\infty$.

The Kodaira--Spencer cocycle of a nonzero $t$-variation is a
nonzero constant times $\zeta\partial_\zeta$ on this annulus.
The cotangent space to the labelled four-point moduli at this
sphere is one-dimensional, with basis
$Q=dw^2/((w-a)(w-b)(w-c))$. This follows either from Serre
duality or from $K_{\PP^1}^{\otimes2}(a+b+c+\infty)$ having
degree zero. The cocycle pairs with $Q$ by the contour integral
of
\[
 \frac{\zeta\phi'(\zeta)^2\,d\zeta}
      {(\phi(\zeta)-a)(\phi(\zeta)-b)(\phi(\zeta)-c)}.
\]
The only poles inside the unit circle are $A,B$. The sum of
their residues is
\[
 \frac{2}{(A-B)(A+B-2R)}
 \left[
 \frac{A(A-R)}{(A-C)(A+C-2R)}
 -
 \frac{B(B-R)}{(B-C)(B+C-2R)}
 \right]=0.
\]
Indeed solving this cancellation, after cancelling $A-B$,
gives $A=C(R-B)(C-2R)/(C^2-2RC+RB)=420/569$.
Thus the pairing with the whole cotangent space vanishes,
and the tangent map is zero. The family is a local chart
near its nodal parameter by Proposition~\ref{A:plumbing},
so this is a nonconstant family with a critical point.
The exact-arithmetic program
\texttt{code/A\_plumbing\_certificate.py} verifies all ten
finite parameter and residue checks, including the locations
of all six distinct poles. This example does not contradict
the local theorem near $t=0$ or the everywhere invertible
real radius map, which uses moving canonical cap coordinates.
% Suggested insertion in the proof of A:hinich-explicit.
% Independent audit by /root/algebra_c, 19 November2026.
\paragraph{An independent explicit cross-ratio verification.}
For these particular coordinates one can compute the entire plumbing family directly, without using the cotangent pairing. The first capped component has a global coordinate $\zeta_1$: on the original bounded component it is the inverse of the injective map $w=\phi(\zeta_1)$, and it is the exterior coordinate on its attached cap. Its three marks are $A,B,\infty$. The second capped component has the original global $w$-coordinate: on its attached interior cap use $w=\phi(\zeta_2)$. Its three marks are $c,\infty,R^2$. The two node coordinates are therefore exactly
\[
 x=(1-\varepsilon)/\zeta_1,
 \qquad y=\zeta_2/(1+\varepsilon).
\]

Fix any $t$ with $0<|t|<1$, and put $\alpha=t/t_0$. The first retained plumbing piece is
\[
 |\zeta_1|<(1-\varepsilon)/|t|,
\]
and the second retained piece is the complement in the $w$-sphere of
\[
 \phi\bigl(\{|\zeta_2|\le(1+\varepsilon)|t|\}\bigr).
\]
On the first piece define $w_t=\phi(\alpha\zeta_1)$; on the second use its existing $w$-coordinate. The plumbing equation $xy=t$ is equivalent to $\zeta_2=\alpha\zeta_1$, so these functions agree on the overlap. Because
\[
 |\alpha|\,(1-\varepsilon)/|t|
       =(1-\varepsilon)/t_0=1+\varepsilon<R,
\]
and $\phi$ is injective on $|\zeta|<1+\varepsilon$, the first piece maps biholomorphically to $\phi(\{|\zeta|<1+\varepsilon\})$. Its overlap with the second piece is exactly the image of
\[
 (1+\varepsilon)|t|<|\zeta_2|<1+\varepsilon.
\]
The two images cover the $w$-sphere. Thus $w_t$ is a global conformal coordinate on the plumbed sphere for every $0<|t|<1$, not merely near $t_0$.

The four labelled marks in this coordinate are
\[
 a(t)=\phi((t/t_0)A),\quad
 b(t)=\phi((t/t_0)B),\quad c=\phi(C),\quad\infty.
\]
They stay distinct throughout this punctured unit disc: the first two lie in the injective $\phi$-disc, while $c,\infty$ are outside it. A cross-ratio coordinate on labelled $M_{0,4}$ is therefore the explicit rational function
\[
 \Lambda(t)=\frac{\phi((t/t_0)A)-c}{\phi((t/t_0)B)-c}.
\]
Substituting the specified rational constants and differentiating gives
\[
 \Lambda'(t)=\frac{201579840}{323761}\,
 \frac{(101t-99)(101t+12518)}
 {(101t-1320)^2(101t+924)^2}.
\]
In particular
\begin{align*}
 \Lambda(t_0)&=\frac{323680}{323761},
 &\Lambda'(t_0)&=0,\\
 \Lambda'(0)&=-\frac{909}{1752520}\ne0,
 &\Lambda''(t_0)&=\frac{22850240}{44933817907}\ne0.
\end{align*}
Thus the nonzero parameter $t_0$ is a simple critical point, while the family is a local smoothing coordinate at the nodal parameter $0$. This verifies the geometric conclusion independently of the Kodaira--Spencer computation. The exact certificate \texttt{code/C\_plumbing\_crossratio.py}, with output \texttt{logs/C\_plumbing\_crossratio.json}, also checks that all zeros and poles of $\Lambda$ lie outside the closed unit disc, and that $\Lambda=1$ occurs there only at $t=0$. Its eleven assertions use exact rational functions.

\end{proof}


\section{A.2. Capping, radii, limits, and collars}

\subsection{A.2(a): a cap cannot depend only on its own length}

\paragraph{Construction and source comparison.}
Here a cap is the embedded pointed disc obtained by attaching a flat
semi-infinite cylinder along a geodesic boundary by arclength, then filling
its puncture. Its boundary is the original geodesic seam. This is the
infinite-grafting construction in Mondello (2011), Theorem~5.4
(the downloaded source has the label \texttt{thm:grafting}), restricted
to the smooth finite stratum. F\i rat (2021), equations~(1.2), (1.3a),
and the projective-connection and cap-coordinate formulas in \S3,
describes the same construction
on a pair of pants: the boundary lengths are $L_i=2\pi\lambda_i$,
the cap coordinates are $w_i$, and the flat metric on a cap is
$\lambda_i^2|dw_i|^2/|w_i|^2$. We opened both sources and their
downloaded \LaTeX{} files. We now give a direct argument that does not
require choosing branches of the gamma functions in the latter paper.

\begin{theorem}[Proved: a single cap remembers all the lengths]
\label{A:single-cap-all-lengths}
Let $(X;p_1,\ldots,p_n)$ be a fixed connected compact marked Riemann
surface. Suppose two systems of pairwise disjoint embedded pointed
discs $(D_j,p_j)$ and $(D'_j,p_j)$ arise by infinite grafting of
hyperbolic surfaces with geodesic boundaries of positive lengths
$\ell_j$ and $\ell'_j$, respectively. If, for one index $i$,
\[
 D_i=D'_i,
\]
then $\ell_j=\ell'_j$ for every $j$, including $j=i$.
With the usual marking convention,
Mondello's grafting theorem further implies that the two complete
systems are equal.
\end{theorem}

\begin{proof}
A hyperbolic surface carries its Fuchsian complex projective structure:
on its universal cover a local isometry into the upper half-plane is a
holomorphic developing map and two choices differ by an element of
$\operatorname{PSL}_2(\mathbb R)$. Near a lift of a geodesic boundary
of length $\ell$, choose a logarithmic developing coordinate
$\zeta=t+is$ such that the boundary is $s=0$, its arclength coordinate
is $t$, its deck transformation is $\zeta\mapsto\zeta+\ell$, and
the developing map is $i e^{\zeta}$ up to a M\"obius transformation.
The direction of $s$ can be chosen to point into the attached cylinder;
reversing this convention only reverses the sign of the exponent below.
The grafted cylinder has exactly the same complex coordinate $\zeta$.
Consequently $i e^{\zeta}$ extends the projective structure across the
seam and throughout the cylinder. It is locally univalent everywhere
on the cylinder. Its monodromy under one turn is multiplication by
$e^{\ell}$, a projective transformation, so this continuation defines
a projective structure on the punctured capped surface.

On the filled cap let
\[
 w=\exp(2\pi i\zeta/\ell),\qquad
 \lambda=\ell/(2\pi).
\]
Thus $w$ maps the cap to the unit disc, takes its centre to $0$, and
is determined by the pointed embedded disc up to multiplication by a
constant of absolute value~1. The developing map on any simply
connected subregion of the punctured cap is, up to a projective change
of coordinate, $w^{-i\lambda}$. A different sign convention gives
$w^{i\lambda}$ and the same Schwarzian derivative.

Write $\{f,z\}=f'''/f'-\frac32(f''/f')^2$. Differentiating
$w^\alpha$ gives
\[
 \{w^\alpha,w\}=\frac{1-\alpha^2}{2w^2}.
\]
The Schwarzian chain rule therefore gives the projective connection
in any holomorphic coordinate $z$ on the cap as
\begin{equation}\label{A:cap-projective-connection}
 Q(z)=\{w,z\}+
 \frac{1+\lambda^2}{2}\left(\frac{w'(z)}{w(z)}\right)^2.
\end{equation}
This expression is invariant under the allowed rotations of $w$.
Since $w$ has a simple zero at the centre and $w'$ does not vanish,
$\{w,z\}$ is holomorphic there, whereas the second term has a pole
of order two with leading coefficient $(1+\lambda^2)/2$. We have
thus proved, rather than assumed, that the connection extends to a
meromorphic projective connection on $X$ with regular singularities
at its marked points and with precisely these leading coefficients.
There are no other poles: the developing maps on the hyperbolic core,
the collars and the cylinders are locally univalent.

Apply this construction to the two cap systems. We first prove the
stronger assertion that their common cap determines its own length.
Suppose $D_i=D'_i$ but $\lambda_i\ne\lambda'_i$, and use their
common normalized coordinate $w$, unique up to rotation. Since
$Q-Q'$ is a globally defined meromorphic quadratic differential,
so is
\[
 \Phi=\frac{Q-Q'}{(\lambda_i^2-(\lambda'_i)^2)/2}.
\]
On the cap, $\Phi=(dw/w)^2$. Choose a developing map $u$ of the
first projective structure, with $u=w^{i\lambda_i}$ on a simply
connected portion of this cap. Then
\[
 (du/u)^2=-\lambda_i^2\Phi
\]
there, hence on the whole universal cover of the punctured surface
by meromorphic continuation. If $M$ is any projective monodromy
transformation, deck invariance of the right side implies
\[
 \left(\frac{d(M(u))}{M(u)}\right)^2=(du/u)^2.
\]
The developing image contains an open set, so the rational function
identity $(M'(z)/M(z))^2=z^{-2}$ follows. On a connected open set
its square root has a constant sign, giving
$M'(z)/M(z)=1/z$ or $-1/z$. Integration, and then rational
continuation, gives respectively $M(z)=az$ or $M(z)=a/z$.
Thus all holonomy lies in the normalizer of $\{0,\infty\}$,
which has an abelian subgroup of index two.

This is impossible. Attaching punctured cylinders is a homotopy
equivalence, so this holonomy contains the faithful Fuchsian
holonomy of the hyperbolic core. A finite-area hyperbolic core with
geodesic boundary has $\chi<0$ by Gauss--Bonnet. With a nonempty
boundary its fundamental group is free of rank
$2g+n-1\ge2$. Such a group is not virtually abelian: an index-$d$
subgroup is the fundamental group of a $d$-sheeted covering of a
rank-$r$ graph, and has rank $1+d(r-1)\ge2$. This contradiction
proves $\lambda_i=\lambda'_i$.

On the common cap $D_i$, the two normalized disc coordinates differ
by a rotation.
Equation~\eqref{A:cap-projective-connection} and
$\ell_i=\ell'_i$ show that the two projective connections coincide
there. Their difference is a globally defined meromorphic quadratic
differential on $X$: the Schwarzian coordinate-change terms cancel.
It vanishes on a nonempty open subset, hence vanishes everywhere by
the identity theorem. Comparing its leading double-pole coefficients
at each $p_j$ gives
\[
 1+(\ell_j/2\pi)^2=1+(\ell'_j/2\pi)^2.
\]
All lengths are positive, so $\ell_j=\ell'_j$ for every $j$.
Finally the restriction of Mondello's map
$(\mathrm{gr}_\infty,L)$ to positive finite lengths is injective;
the fixed capped marked surface and the common length vector therefore
determine the same marked bordered surface and its attached caps.
This last invocation is needed only for equality of the whole systems,
not for the equality of their lengths just proved.
\end{proof}

\begin{corollary}[Proved extension: positive caps distinguish cuspidal ends]
\label{A:single-cap-with-cusps}
The preceding uniqueness holds also when some of the other ends are
complete hyperbolic cusps, assigned length zero, while the compared
cap is a positive cap in both constructions. Therefore keeping one
positive pointed cap unchanged cannot be compatible with replacing
any other positive geodesic boundary by a complete cusp on the same
marked compactification.
\end{corollary}

\begin{proof}
At a complete cusp a parabolic developing coordinate is, up to a
projective transformation and a nonzero scalar, $\log w$, where
$w$ is a holomorphic cusp coordinate vanishing simply at the filled
point. Direct differentiation gives
$\{\log w,w\}=1/(2w^2)$. Hence the projective connection extends
meromorphically there with the same leading coefficient
$(1+\lambda^2)/2$ at $\lambda=0$. The compared cap is positive,
so the holonomy argument above applies unchanged and first recovers
its own positive length. The identity theorem then identifies the
two global connections. Their double-pole coefficients recover all
the nonnegative lengths, including which ones are zero. In
particular a positive length cannot become zero.
\end{proof}

\begin{proposition}[Refuted under the literal cusp admissibility convention]
\label{A:MDB-forgetting-admissibility}
Interpret the admissibility of the external cap system in
67/075(b) to mean that on each normalization component the node
branches are complete cusps, as the point-like boundary components
$s,t$ on 67/081 specify. Under this convention, the proposed
capping-and-forgetting operation
$MDB_{g,n}\to M_{\mathrm{st}}ND_{g,n}$ need not take values in
its stated target. It already fails for two positive-boundary
pairs of pants glued along one boundary, giving type $(0,4)$.
More generally it fails on any normalization component that has
both a forgotten positive node cap and a retained positive
external cap.
\end{proposition}

\begin{proof}
Choose two hyperbolic pairs of pants with three positive boundary
lengths each, the two boundaries to be glued having equal lengths.
Glue these geodesic boundaries by an arclength-reversing isometry.
This produces a compact bordered hyperbolic four-holed sphere with
a specified separating geodesic; hence it is a legitimate source
object. Cut it again and attach its canonical infinite cylinders.
The compactified components are three-pointed spheres, each
carrying a full system of three positive canonical caps. Forget
the cap at the branch corresponding to the separating geodesic,
but keep the two external caps and the node's nonzero tangent
tensor, exactly as prescribed on 67/075.

If the remaining caps were admissible with the node branch treated
as a cusp, the same compactified three-pointed sphere would be
obtained by grafting a second hyperbolic surface with that end of
length zero and with the retained external caps unchanged.
Choose either retained positive cap. Corollary~\ref{A:single-cap-with-cusps}
forces equality of all three lengths in the two constructions,
contradicting positivity of the first construction's node-boundary
length and zero length of the second. The same argument is local
to a normalization component and proves the general assertion.
\end{proof}

\paragraph{The two meanings of a partial admissible system.}
This is a conditional interpretation result, not an assertion that
the manuscript formally defines its target category. The wording
on 67/075(b) requires an admissible system only around external
marks, and the 67/081 pants target explicitly changes the glued
boundaries into point-like components. We checked the latter phrase
against the original scan. A different, precise definition is to
call a partial cap system \emph{extendible} when it occurs by
forgetting selected caps from a full admissible system. That
definition makes the displayed forgetting operation well-defined,
but it does not identify the retained caps with the canonical caps
obtained after declaring omitted ends to be cusps. With the nodal
tensor still independent and unrestricted, the radius obstruction
of Theorem~\ref{A:04} remains: even an extendible external system
cannot produce a node tensor of norm $17$ in its fixed normalized
three-pointed coordinates, since every possible node cap has
conformal radius at most $4$ and their product at most $16$.

\begin{corollary}[Refuted: the question marked ``?'' at 67/068]
\label{A:pants-cap-dependence}
Normalize the capped, labelled pair of pants to
$(\mathbb P^1;0,1,\infty)$. The cap at $0$ is not a function of
its own boundary length alone. Specifically the length vectors
\[
 (2\pi,2\pi,2\pi),\qquad(2\pi,4\pi,2\pi)
\]
give different embedded pointed discs at $0$, although that boundary
has the same length in both examples. Every positive length vector is
realized by a hyperbolic pair of pants, so both examples meet all
hypotheses of the question.
\end{corollary}

\begin{proof}
The existence statement follows either from the usual gluing of two
right-angled hyperbolic hexagons with alternating side lengths
$\ell_1/2,\ell_2/2,\ell_3/2$, or directly from the positive-length
stratum of the opened grafting theorem. The three filled punctures
have the unique labelled normalization $0,1,\infty$. If the two caps
at $0$ were the same, Theorem~\ref{A:single-cap-all-lengths} would
force the second boundary lengths to be equal, a contradiction.

For an explicit local witness, put
$\Delta_j=(1+\lambda_j^2)/2$. A projective connection on the
three-punctured sphere with only these regular singularities must be
\begin{equation}\label{A:pants-projective-connection}
 Q(z)=\frac{\Delta_1}{z^2}
      +\frac{\Delta_2}{(z-1)^2}
      +\frac{\Delta_3-\Delta_1-\Delta_2}{z(z-1)}.
\end{equation}
Indeed the difference of two such connections has at most simple
poles at the three marked points, and would be a section of
$K_{\mathbb P^1}^{\otimes2}(0+1+\infty)$, of degree $-1$;
there is no nonzero such section. Alternatively partial fractions
and the coordinate $1/z$ verify the displayed formula directly.
Let $w(z)=c(z+a z^2+O(z^3))$ be the unit-disc coordinate on the
cap at $0$, where $c\ne0$. Its Schwarzian is holomorphic at $0$,
and $(w'/w)^2=z^{-2}+2a z^{-1}+O(1)$.
Comparison of the simple-pole coefficients in
\eqref{A:cap-projective-connection} and
\eqref{A:pants-projective-connection} yields
\[
 a=\frac{\Delta_1+\Delta_2-\Delta_3}{2\Delta_1}
  =\frac{1+\lambda_1^2+\lambda_2^2-\lambda_3^2}
         {2(1+\lambda_1^2)}.
\]
The two examples give $a=1/2$ and $a=5/4$. The ratio of the
quadratic and linear coefficients is unchanged by rotating $w$;
equal pointed discs would have equal ratios. This is also exactly
the independently opened expansion in F\i rat (2021), equation
(1.3a). The program \texttt{code/A567\_matrices.py} computes these
two rational numbers exactly, and its JSON output records the result.
\end{proof}

\paragraph{Scope.}
This proves the stated dependence on the other lengths; a single
positive cap in fact determines the entire length vector. It does
not by itself establish that
the vector of conformal radii determines the cap system: a conformal
radius discards the higher Taylor coefficients used in this proof.
Nor does the argument identify a one-variable limiting cap when several
lengths vary without a specified direction.


\begin{theorem}[Canonical capping is infinite grafting]\label{A:grafting}
\status{Known, with a proved identification of the manuscript operation.}
Let $S$ be a compact oriented surface with $n\geq1$ labelled boundary
components and $\chi(S)=2-2g-n<0$. Let $\mathcal T^B(S)$ denote marked
hyperbolic surfaces with these boundary components geodesic, and let
$\mathcal T_{g,n}$ denote the Teichmuller space of the corresponding
punctured surfaces. Canonical capping of 67/061--070, followed by
remembering the cap centres, satisfies
\[
 \Psi=(\operatorname{cap},\ell):\mathcal T^B(S)
 \xrightarrow{\ \sim\ }\mathcal T_{g,n}\times(0,\infty)^n
\]
as a mapping-class-equivariant real analytic diffeomorphism. Allowing
some boundary lengths to vanish extends this to a homeomorphism
\[
 \overline\Psi:\mathcal T^{B,\geq0}(S)
 \xrightarrow{\ \sim\ }\mathcal T_{g,n}\times[0,\infty)^n,
\]
real analytic on every fixed boundary/cusp stratum. No assertion of
analyticity across all corner strata is needed here.
\end{theorem}
\begin{proof}
On a boundary of length $\ell>0$, write arclength as
$s\in\R/\ell\Z$. The semi-infinite Euclidean cylinder with coordinates
$(s,t)\in(\R/\ell\Z)\times[0,\infty)$ is identified conformally with
the punctured closed disc by
\[
 \zeta=\exp\!\left(\frac{2\pi i s-2\pi t}{\ell}\right).
\]
Its boundary circle has exactly the rotational torsor given by
arclength. Adding $\zeta=0$ is precisely the manuscript's pointed
conformal cap. Translating the arclength origin rotates $\zeta$,
which does not change the pointed cap. Thus the punctured capped
surface is Mondello's $\operatorname{gr}_\infty(S)$, with identical
marking and boundary lengths. Mondello's Theorem~5.4 \citep[p.~660]{C_Mondello2011} now gives the
stated maps and regularity, after restricting his finite part to
smooth surfaces and then to the required boundary/cusp strata.
His definitions 2.1--2.3 allow precisely these collapsed boundary
components. This proves that all hypotheses of the cited theorem
hold for the manuscript operation. It does not identify lengths with
conformal shadows.
\end{proof}

% Independent contribution to Part A.2(b), by /root/algebra_c.
% AI-produced by GPT-6 Astra (OpenAI), ChatGPT Pro mode, 19 November 2026.
\subsection{A positive energy identity for the conformal-radius map}
\label{C:A2b-energy-section}

Fix a compact connected Riemann surface $X$ of genus $g$, distinct labelled points $p_1,\ldots,p_n$, and local holomorphic coordinates $z_i$ with $z_i(p_i)=0$. Assume $2g-2+n>0$. The coordinates specify the norms in which we measure conformal radii. For boundary lengths $\ell_i\geq0$, Mondello's infinite-grafting theorem, Theorem~5.4 (2011, p.~660), gives the unique marked hyperbolic surface whose infinite grafting is the fixed pointed conformal surface $X$ and whose boundary lengths are $\ell_i$~\cite{C_Mondello2011}. Set $\lambda_i=\ell_i/(2\pi)$. For $\lambda_i>0$, denote its canonical cap by $D_i(\lambda)$ and its conformal radius in $z_i$ by $r_i(\lambda)>0$. For $\lambda_i=0$, the cap is the marked point and we set $r_i=0$.

We first work on a fixed face: the active index set $I=\{i:\lambda_i>0\}$ is fixed, and all other ends remain cusps. Put $x_i=\lambda_i^2$ for $i\in I$, and define
\[
 F_X:(0,\infty)^I\longrightarrow\mathbb R^I,
 \qquad F_X(x)_i=\log r_i(\sqrt{x}).
\]
The following argument concerns every stable $(g,n)$, including $(0,3)$, $(0,4)$ and $(1,1)$.

\begin{lemma}[Known regularity, with its application spelled out]
\label{C:A2b-regularity}
The caps and their normalized conformal coordinates vary real-analytically along each finite-dimensional real-analytic family in a fixed face of the length space. In particular $F_X$ is real-analytic. After identifying the infinitely grafted surface with $X\setminus\{p_i\}$, its Thurston metric varies differentiably off the moving seams. At a seam it is $C^{1,1}$, its first variation is continuous, and the one-sided derivatives used below exist.
\end{lemma}
\begin{proof}
Mondello's Theorem~5.4 states real-analyticity of the inverse on each finite topological stratum. We explain why it provides the local data needed here. A family of the resulting hyperbolic surfaces can be represented by varying Fenchel--Nielsen data; its boundary geodesics and Fermi collars vary real-analytically. With $s$ the signed hyperbolic distance from a boundary geodesic and $\theta$ its angle of period $2\pi$, the collar metric is
\[
 ds^2+\lambda_i^2\cosh^2(s)\,d\theta^2.
\]
The change $dy=ds/\cosh s$, $u=\lambda_i\theta$ is conformal and turns this into $\sec^2(y)(du^2+dy^2)$, with the seam at $y=0$. Continuing with $du^2+dy^2$ on $y\geq0$ gives the flat cylinder. For identifications with a fixed differentiable surface, keep $\theta$ and $y/\lambda_i$ fixed on each end and in a small collar. The metrics there change by a scalar conformal factor, so the Beltrami coefficient is zero on fixed neighborhoods of the punctures and seams. Interpolate these identifications over the remaining compact core. Its Beltrami coefficient depends real-analytically in $L^\infty$, while the seams stay fixed in this preliminary model. This is the local family construction underlying Mondello's proof in \S5.5, pp.~661--663.

Normalized solutions of the Beltrami equation depend differentiably, and real-analytically for analytic data, on these parameters: Ahlfors--Bers (1960), Theorem~10, Theorem~11 and its corollary, pp.~402--404~\cite{C_AhlforsBers1960}. Apply this to the marked conformal identifications with the fixed $X$. One may perform the construction in a finite conformal atlas, or on its normalized universal cover; uniqueness of the normalized solution makes the identifications compatible on overlaps. Because the conformal class is fixed, their targets are the same marked $X$. The preliminary end identifications just chosen are conformal, so the normalized Beltrami solutions are holomorphic on their fixed puncture and seam neighborhoods. Uniform parameter differentiability there implies differentiability of every holomorphic derivative by Cauchy's formula. On a cylinder the exponential
$w_i=\exp((-y+iu)/\lambda_i)$ is holomorphic and extends through the seam by the collar formula. Its normalization at $p_i$ fixes the harmless rotation. The parameter derivatives in the interior are holomorphic, and Cauchy's formula gives their derivatives on compact subsets. The same argument applies in the extended collar, so it includes both one-sided seam derivatives. Finally the two metric factors $1$ and $\sec^2(y)$ agree to first order at $y=0$ and have bounded second derivatives there. This gives the stated $C^{1,1}$ regularity and continuity of the first variation. Thus the cited parameter-dependence statements apply to the actual cap coordinates, not merely to abstract moduli points.
\end{proof}

For a tangent vector $v=(\dot x_i)_{i\in I}$, take a differentiable curve $x(t)$ with $x(0)=x$ and $\dot x(0)=v$. Write $g_t$ for its Thurston metric on the fixed punctured conformal surface, and let
\[
 \psi_v=\left.\frac{d}{dt}\log(g_t/g_0)\right|_{t=0}.
\]
This is a scalar function, independent of the holomorphic coordinate used to express the two conformal metrics. The base hyperbolic core is denoted $M$, and the base flat end belonging to $i$ by $Z_i$.

\begin{theorem}[Proved: the radius energy identity]
\label{C:A2b-energy}
For any two tangent vectors $v,w\in\mathbb R^I$, there are real harmonic functions $V_{i,v},V_{i,w}$ on the flat cylinders $Z_i$ such that
\begin{align}
4\pi\sum_{i\in I}v_i\,(dF_X)_x(w)_i
={}&\int_{X\setminus\{p_i\}}
 \bigl(\langle\nabla\psi_v,\nabla\psi_w\rangle_{g_0}
       +2\,1_M\psi_v\psi_w\bigr)\,dA_{g_0}\notag\\
 &+4\sum_{i\in I}\int_{Z_i}
       \langle\nabla V_{i,v},\nabla V_{i,w}\rangle\,dA.
\label{C:eq-radius-energy}
\end{align}
Every integral converges. The right-hand side is symmetric in $v,w$ and strictly positive when $v=w\ne0$. Consequently $DF_X(x)$ is symmetric positive definite at every $x$.
\end{theorem}
\begin{proof}
We give all normalizations and signs. Fix an active end and abbreviate its base length parameter by $\lambda$. Use flat coordinates
\[
 (u,y)\in(\mathbb R/(2\pi\lambda\mathbb Z))\times[0,\infty),
 \qquad w_0=\exp((-y+iu)/\lambda).
\]
The coordinate $w_t$ on the moving cap satisfies $|w_t|=1$ on its seam and has a simple zero at the fixed puncture. By Lemma~\ref{C:A2b-regularity}, it extends across the seam for small $t$. Its logarithmic variation
\[
 h_v=\left.\frac{d}{dt}\log w_t\right|_0
       =\frac{\dot w_t|_0}{w_0}
\]
is a single-valued holomorphic function on the base cap including its center and a collar of its boundary. Indeed $w_t$ vanishes at the fixed center, so the quotient has a removable singularity there. Define
\[
 a_v=\dot\lambda/\lambda,\qquad V_v=\lambda\operatorname{Re}h_v.
\]
Then $V_v$ is harmonic on $Z_i$. In the chosen local coordinate at the puncture,
$w_t=e^{i\alpha(t)}z_i/r_i(t)+O(z_i^2)$, and hence
\begin{equation}
 \lim_{y\to\infty}V_v(u,y)
   =-\lambda\,d(\log r_i)(v)=:b_{i,v}.
 \label{C:eq-radius-end}
\end{equation}
The convergence, with all derivatives, is exponential in $y/\lambda$. In particular the mean of $V_v$ on any horizontal circle is the constant $b_{i,v}$.

On the flat cap the metric is $g_t=\lambda(t)^2|dw_t/w_t|^2$. Differentiating this identity in any fixed holomorphic coordinate gives
\[
 \psi_v=2a_v+2\operatorname{Re}(w_0\partial_{w_0}h_v)
        =2a_v-2\partial_y V_v.
\]
Differentiating $\log|w_t|=0$ at its moving seam shows that the seam's normal velocity, positive from the hyperbolic core into the flat cylinder, is exactly $V_v(u,0)$. The base coordinate $u+iy$ extends to a hyperbolic collar $y<0$, where the metric is $\sec^2(y)|d(u+iy)|^2$. Thus the logarithmic metric factor and its first derivative agree across $y=0$, while the difference between its hyperbolic and flat second normal derivatives is $2$.

The curvature equation now gives, in the interiors,
\[
 \Delta_{g_0}\psi_v=2\psi_v\quad\hbox{on }M,
 \qquad \Delta_{g_0}\psi_v=0\quad\hbox{on each }Z_i,
\]
where $\Delta$ is the positive-coordinate Laplacian $\partial_u^2+\partial_y^2$ in a flat chart. Its interface condition is
\begin{equation}
 (\partial_y\psi_v)_{+}-(\partial_y\psi_v)_{-}=2V_v.
 \label{C:eq-radius-jump}
\end{equation}
One can obtain this by differentiating the matched first normal derivatives. Equivalently, the core is locally $\{y<tV_v(u,0)+o(t)\}$, so differentiating
$\Delta\phi_t=2e^{\phi_t}1_{M_t}$ gives the distributional term $2V_v\delta_{y=0}$. This also fixes the sign in~\eqref{C:eq-radius-jump}.

For clarity, all integrations can first be performed after truncating the flat cylinders at $y=T$. On a fixed cuspidal end, use its normalized local cusp coordinate $q_t$. The formula
\[
 g_t=\frac{|dq_t/q_t|^2}{(\log|q_t|)^2}
\]
and holomorphy of $\dot q_t/q_0$ show that, in $Y=-\log|q_0|$,
$\psi_v=O(Y^{-1})$ and $\partial_Y\psi_v=O(Y^{-2})$, up to exponentially decaying terms. These estimates make the cusp Dirichlet and potential integrals finite, and make its boundary terms tend to zero. On flat ends, $\psi_v\to2a_v$ and its derivatives decay exponentially; the corresponding boundary terms also tend to zero. These observations justify every following integration by parts and the passage to the untruncated surface.

Apply Green's identity to $\psi_v$ and the equation for $\psi_w$ on the core and all cylinders separately. Their outward normals at the seam are respectively $+\partial_y$ and $-\partial_y$. Using continuity of $\psi_v$ and~\eqref{C:eq-radius-jump}, we obtain
\begin{align}
 E_\psi(v,w)&:=\int
 \bigl(\langle\nabla\psi_v,\nabla\psi_w\rangle
          +2\,1_M\psi_v\psi_w\bigr)\,dA\notag\\
 &=-2\sum_{i\in I}\int_{\partial M_i}\psi_v V_{i,w}\,du.
 \label{C:eq-radius-green}
\end{align}
Harmonicity of $V_{i,v}$ similarly gives
\[
 E_{V_i}(v,w):=\int_{Z_i}\langle\nabla V_{i,v},\nabla V_{i,w}\rangle\,dA
 =-\int_{\partial M_i}(\partial_y V_{i,v})V_{i,w}\,du.
\]
Substitute $\psi_v=2a_{i,v}-2\partial_yV_{i,v}$ into~\eqref{C:eq-radius-green}. Since the circle has length $2\pi\lambda_i$ and the mean of $V_{i,w}$ is~\eqref{C:eq-radius-end}, this yields
\begin{align*}
 E_\psi(v,w)
 &=-4\sum_i a_{i,v}(2\pi\lambda_i)
                  (-\lambda_i\,d\log r_i(w))
       -4\sum_iE_{V_i}(v,w)\\
 &=4\pi\sum_i d(\lambda_i^2)(v)\,d\log r_i(w)
       -4\sum_iE_{V_i}(v,w).
\end{align*}
This is~\eqref{C:eq-radius-energy}.

The energy forms on its right are symmetric. If $v=w$ and the right-hand side vanishes, then $\psi_v=0$ on the nonempty hyperbolic core and $\nabla\psi_v=0$ everywhere. Each $V_{i,v}$ has zero Dirichlet energy and is constant. Continuity at the seam and $\psi_v=2a_{i,v}-2\partial_yV_{i,v}$ then give $a_{i,v}=0$ for every active $i$. Hence $v_i=2\lambda_i\dot\lambda_i=0$ for every $i$. Conversely zero variation gives zero energy. This proves strict positive definiteness.
\end{proof}

\begin{corollary}[Proved: global injectivity and a strictly convex potential]
\label{C:A2b-injectivity}
On every face of the nonnegative boundary-length space, the canonical cap system is determined by its conformal radii. On its positive coordinates the logarithmic radius map is a real-analytic diffeomorphism onto an open subset of $\mathbb R^I$, and is the gradient of a strictly convex real-analytic function of $x_i=(\ell_i/2\pi)^2$. In particular the full map of admissible canonical cap systems to their tangent-disc radii is injective. Changing the local coordinate norms translates the logarithmic radius map by a constant vector and does not change these assertions.
\end{corollary}
\begin{proof}
For distinct $x,x'$ in the convex positive orthant, put $v=x'-x$. The fundamental theorem of calculus and Theorem~\ref{C:A2b-energy} give
\[
 v\cdot(F_X(x')-F_X(x))
 =\int_0^1v^TDF_X(x+tv)v\,dt>0.
\]
Thus $F_X$ is injective. Its positive definite derivative is invertible, so the inverse function theorem makes it a local real-analytic diffeomorphism; injectivity then glues its local inverses into the claimed inverse on the open image. Symmetry of its derivative says that the one-form $\sum_i(F_X)_i\,dx_i$ is closed. Since the positive orthant is convex, it has a primitive $S_X$, whose Hessian is $DF_X$ and hence is positive definite. Explicitly, fixing $x_*$ one can take
\[
 S_X(x)=\int_0^1F_X(x_*+t(x-x_*))\cdot(x-x_*)\,dt.
\]
Its additive constant has no geometric significance.

If two radius vectors in the full nonnegative parameter space agree, their zero entries specify the same cuspidal indices; the remaining entries lie in the same face, where the preceding argument identifies their lengths. Mondello's uniqueness then identifies their cap systems on the fixed marked $X$. The all-cuspidal face consists of the single complete hyperbolic surface and requires no argument. Finally a coordinate change multiplies each radius by its fixed derivative modulus at $p_i$, so it adds a constant to $F_X$ and changes neither its derivative nor injectivity.
\end{proof}

\paragraph{The boundary argument needed for the image.}
The energy identity proves injectivity and local coordinates. To identify the open image, the next subsection supplies a separate support-gap estimate and properness argument. The image description is not inferred from positive definiteness alone.

% Full image and limit contribution to A.2(b,c), by /root/algebra_c.
% AI-produced by GPT-6 Astra (OpenAI), ChatGPT Pro mode, 19 November 2026.
\subsection{The exact image: simultaneous weighted conformal-radius inequalities}
\label{C:A2b-image-section}

Continue with a fixed stable pointed conformal surface $(X,p_1,\ldots,p_n)$ and its specified tangent coordinate norms. Let $\mathcal D_X$ be the systems of pairwise disjoint open conformal discs $D_i\ni p_i$, each avoiding every other marked point. Their boundaries need not be disjoint. Write $r_i(D)$ for their conformal radii and define, for $a\in\mathbb R_{\geq0}^n$,
\[
 H_X(a)=\sup_{D\in\mathcal D_X}\sum_i a_i\log r_i(D),\qquad H_X(0)=0.
\]
This definition is intrinsic once the tangent norms are fixed. The next statement computes the support function by a unique quadratic differential, so the eventual image test does not require a search over arbitrary discs.

\begin{theorem}[Known extremal theorem, with the complexity-zero case proved]
\label{C:A2b-support}
For $a\ne0$, let $Q_a$ be the unique Jenkins--Strebel differential whose characteristic domains are exactly the punctured discs at the active punctures $\{p_i:a_i>0\}$, with circumferences $2\pi\sqrt{a_i}$; at the inactive marked points it has at most simple poles. In particular no additional characteristic annuli are permitted. Let $D_i^a$ be its characteristic pointed discs and $R_i(a)$ their conformal radii. Then
\[
 H_X(a)=\sum_{a_i>0}a_i\log R_i(a)<\infty.
\]
Among disjoint active pointed discs avoiding all the other marked points, equality in this weighted radius inequality holds exactly for the indicated Strebel discs. Their complement is the critical graph, including inactive marked points, and has empty interior. The function $H_X$ is convex, positively homogeneous, and continuous on the whole nonnegative orthant.
\end{theorem}
\begin{proof}
For $3g+n-3>0$, existence and uniqueness are Liu (2008), Theorem~3.4; the extremal inequality and its equality clause are Lemma~1.7, pp.~218--219~\cite{C_Liu2008}. Take no characteristic annuli and circumference $2\pi\sqrt{a_i}$ at each active puncture. Lemma~1.2 identifies reduced modulus with $(2\pi)^{-1}\log r$. Substitution in Lemma~1.7 and division by $2\pi$ give exactly the displayed weighted inequality. Liu permits an arbitrary distinguished subset of the punctures, so inactive marked points remain excluded from all candidate discs. The stated graph complement is the characteristic-domain decomposition of this differential. Shrink each active disc by a factor $t<1$ in its Riemann coordinate and insert small discs at the inactive points. Letting $t\uparrow1$ proves that the supremum defining $H_X$ over complete systems has the same value as the active-disc extremum.

Liu assumes $3g+n-3>0$, which excludes the stable case $(0,3)$. We supply the exact reduction. Normalize its points to $0,1,\infty$ and take the double covering
\[
 \pi:\mathbb P^1_w\longrightarrow\mathbb P^1_z,\qquad z=w^2.
\]
The upstairs surface is marked at $0,1,-1,\infty$, so Liu applies. For downstairs weights $(a_0,a_1,a_\infty)$ use upstairs weights
\[
 (b_0,b_1,b_{-1},b_\infty)=(4a_0,a_1,a_1,4a_\infty).
\]
The involution $w\mapsto-w$ preserves these data. Uniqueness makes the upstairs Strebel differential invariant, so it descends to a quadratic differential downstairs. At a branch point invariance forces even Laurent powers, so an upstairs simple pole cannot occur. A regular value or an even-order zero descends to at most a simple pole; a double pole has its coefficient divided by four. Hence the descended differential has exactly the desired active circumferences and no extra singularities. Its characteristic discs descend, giving existence. Conversely the pullback of any downstairs Strebel differential has the specified upstairs characteristic domains; upstairs uniqueness gives downstairs uniqueness.

The same covering transfers the extremal inequality with no loss. A pointed disc at $0$ lifts to a branched degree-two disc with radius $\sqrt{r_0}$, and likewise at $\infty$ in coordinates $1/z,1/w$. A disc at $1$ lifts to two unbranched discs at $\pm1$, each of radius $r_1/2$. For example, if $f(0)=0$, the lift at $0$ is the single-valued univalent map $u\mapsto\sqrt{f(u^2)}$ with the simple zero and chosen derivative. Disjointness is preserved by full preimages. The upstairs weighted sum is therefore
\[
 2\bigl(a_0\log r_0+a_1\log r_1+a_\infty\log r_\infty-a_1\log2\bigr).
\]
The same identity holds for the Strebel radii. The additive constant cancels, proving the downstairs inequality. If equality holds, the upstairs equality clause identifies every active lifted disc, and projection identifies the downstairs discs. Zero weights simply omit the corresponding terms and marked discs from this argument. This covers every stable type.

Convexity and homogeneity follow from the supremum definition. For continuity, let $c_i=H_X(e_i)$ and translate all candidate logarithmic vectors by $-c=( -c_i)$, so their coordinates are nonpositive. Their support function is $H'(a)=H_X(a)-a\cdot c$. It is finite and lower semicontinuous as a supremum of continuous linear functions. If $a^{(m)}\to a\ne0$ in the nonnegative orthant, then for each $\epsilon>0$, eventually $a_i^{(m)}\geq(1-\epsilon)a_i$ for every $i$. For every nonpositive candidate vector $y$, this gives
$a^{(m)}\cdot y\leq(1-\epsilon)a\cdot y$. Taking suprema, then $\limsup$, and then $\epsilon\downarrow0$, yields $\limsup H'(a^{(m)})\leq H'(a)$. Lower semicontinuity gives equality. Continuity at $0$ follows from $H'\leq0$ and from a lower bound furnished by any single fixed small-disc system. Thus $H_X$ is continuous everywhere.
\end{proof}

\begin{proposition}[Proved: direct integral formula for the support data]
\label{C:A2b-support-integral}
The preceding support value can be evaluated directly from $Q_a$. On an active characteristic disc choose the branch
$\sqrt{Q_a}\sim i\sqrt{a_i}\,dz_i/z_i$ near $p_i$, and choose a regular point $b_i$ of its boundary. Then
\[
 w_i(z)=\exp\!\left(-\frac{i}{\sqrt{a_i}}
                  \int_{b_i}^{z}\sqrt{Q_a}\right)
\]
is its normalized disc coordinate up to rotation, and
\[
 \log R_i(a)=\lim_{z\to p_i}
 \left[\log|z_i(z)|-
 \operatorname{Re}\!\left(-\frac{i}{\sqrt{a_i}}
                  \int_{b_i}^{z}\sqrt{Q_a}\right)\right].
\]
Thus the image inequalities use only the unique meromorphic quadratic differential with principal parts $-a_i\,dz_i^2/z_i^2$ whose characteristic domains are precisely the punctured discs around the active marks, followed by these integrals and limits. At inactive marks at most simple poles are allowed; there are no additional characteristic annuli. No elementary expression for its coefficients in every genus is asserted.
\end{proposition}
\begin{proof}
The square root is single-valued on the punctured characteristic disc and has residue $i\sqrt{a_i}$. Its period around the puncture is $-2\pi\sqrt{a_i}$, which becomes $2\pi i$ in the exponential; hence $w_i$ is single-valued. Horizontal boundary arcs make the real part of the exponent zero, so $|w_i|=1$ there. In a cylinder coordinate $\zeta$ one has $Q_a=-a_i(d\zeta/\zeta)^2$ and the chosen square root is $i\sqrt{a_i}\,d\zeta/\zeta$. The displayed expression is therefore $\zeta/\zeta(b_i)$, proving it is a disc coordinate with a simple zero at $p_i$. Since $w_i=e^{i\alpha}z_i/R_i+O(z_i^2)$, taking logarithms gives the limit formula. Changing the boundary basepoint multiplies $w_i$ by a constant of modulus one and leaves $R_i$ unchanged.
\end{proof}

For a nonempty active index set $I\subset\{1,\ldots,n\}$, interpret $a\in\mathbb R_{\geq0}^I$ as zero on the other indices, and define the explicit convex domain
\[
 \mathcal K_{X,I}
 =\left\{y\in\mathbb R^I:
       \sum_{i\in I}a_i y_i<H_X(a)
       \quad\text{for every }a\geq0,\ \sum_{i\in I}a_i=1\right\}.
\]
It is nonempty: subtracting a positive constant from every coordinate of any small-disc logarithmic vector gives a point of this set. It is convex by the defining inequalities. It is open because the positive continuous gap $H_X(a)-a\cdot y$ has a positive minimum on the compact weight simplex.

\begin{proposition}[Proved: uniform support gap]
\label{C:A2b-gap}
For the canonical cap radii $F_X(x)_i=\log r_i$, with active parameters $x_i=\lambda_i^2>0$, we have
\[
 0\leq H_X(x)-x\cdot F_X(x)\leq2g-2+n.
\]
Here inactive ends remain cusps and are included in $n$. The lower inequality is strict for every nonzero finite parameter vector.
\end{proposition}
\begin{proof}
The nonnegative lower bound is Theorem~\ref{C:A2b-support}. For the upper bound, let $g$ be the Thurston metric of the canonical cap system, with hyperbolic core $M$ and boundary lengths $\ell_i=2\pi\lambda_i$. Gauss--Bonnet gives
\[
 A:=\operatorname{Area}_g(M)=2\pi(2g-2+n).
\]
The retraction that sends each flat cylinder $(u,y)$ to $(u,0)$ and is the identity on $M$ is length-nonincreasing. Every loop in the peripheral class of $p_i$ therefore has $g$-length at least $\ell_i$: after retraction it is a loop in the corresponding boundary class of the hyperbolic core. In its hyperbolic double, lifting such a loop joins $z$ to $\gamma_i z$, where $\gamma_i$ has translation length $\ell_i$, so its length is at least $\ell_i$. The fold retraction of the double to the core also shows that the peripheral class has not changed.

Use the Strebel metric $|Q_x|$, whose active cylinders have the same circumferences $\ell_i$. On one of these cylinders write its coordinates as
$(u,h)\in(\mathbb R/\ell_i\mathbb Z)\times(0,\infty)$ and write $g=\rho^2(du^2+dh^2)$. Every horizontal circle is peripheral, so
\[
 \int_0^{\ell_i}\rho(u,h)\,du\geq\ell_i.
\]
Cauchy--Schwarz gives $\int_0^{\ell_i}\rho^2\,du\geq\ell_i$. Integrating in $h$ and summing the cylinders compares the two areas on any union of complete truncated Strebel cylinders.

We state the finite-part calculation to make the cutoff comparison precise. Write $R_i=R_i(x)$ for the Strebel radii. Truncate its $i$th cylinder at
$T_i(\epsilon)=\lambda_i\log(R_i/\epsilon)$. Its flat area is
$\ell_iT_i=2\pi x_i\log(R_i/\epsilon)$. In the fixed local coordinate $z_i$, the removed end is bounded by
$|z_i|=\epsilon+O(\epsilon^2)$, since its normalized Riemann coordinate is $z_i/R_i+O(z_i^2)$. The area of the $O(\epsilon^2)$ discrepancy in either cylindrical metric is $O(\epsilon)$, hence tends to zero. In the canonical metric the corresponding finite part is
\[
 A+2\pi\sum_{i\in I}x_i\log(r_i/\epsilon)+o(1).
\]
Inactive cusps contribute finite area already counted in $A$. The Strebel critical graph has area zero for both metrics; its zeros and the inactive marked points cause no extra area term. Subtracting the common divergent terms from the area inequality and letting $\epsilon\to0$ gives
\[
 A+2\pi\bigl(x\cdot F_X(x)-H_X(x)\bigr)\geq0,
\]
which is the claimed upper bound. Strictness of the lower inequality follows either from the equality clause in Theorem~\ref{C:A2b-support}, since the canonical discs leave an open hyperbolic core, or from the explicit simultaneous collar enlargement proved next.
\end{proof}

\begin{proposition}[Proved: simultaneous enlargement and escape at zero length]
\label{C:A2b-collar}
Put
\[
 Y(\lambda)=\arctan\!\left(\frac{1}{2\sinh(\pi\lambda)}\right),
 \qquad \eta_i=Y(\lambda_i)/\lambda_i.
\]
The canonical cap discs can be simultaneously enlarged to disjoint pointed discs whose logarithmic radii are $F_X(x)_i+\eta_i$. Consequently
\[
 H_X(a)-a\cdot F_X(x)\geq\sum_{i\in I}a_i\eta_i>0
 \quad(a\geq0,\ a\ne0),
\]
and
\[
 F_X(x)_i\leq H_X(e_i)-Y(\lambda_i)/\lambda_i
        \longrightarrow-\infty\quad\text{as }\lambda_i\to0.
\]
The last estimate is uniform in all other lengths.
\end{proposition}
\begin{proof}
Mondello's published Collar Lemma~A.1 (2011, p.~668) gives disjoint embedded half-collars of area $\ell_i/[2\sinh(\ell_i/2)]$~\cite{C_Mondello2011}. In Fermi coordinates their area is $\ell_i\sinh s_i$, so the supplied width is
$s_i=\operatorname{arsinh}(1/[2\sinh(\pi\lambda_i)])$. Its conformal height is
\[
 \int_0^{s_i}\frac{ds}{\cosh s}
 =\arctan(\sinh s_i)=Y(\lambda_i).
\]
The canonical coordinate $w_i$ extends univalently from the flat cap through this collar to $|w_i|<\exp(Y(\lambda_i)/\lambda_i)$. The extended image is a pointed disc and its radius is multiplied by this exponential. The collar lemma gives disjointness and ensures that no other marked point is included. Applying the support inequality to this enlarged system proves the first estimate, and taking $a=e_i$ gives the second. Finally $Y(\lambda)\to\pi/2$ as $\lambda\downarrow0$, proving uniform escape to $-\infty$.
\end{proof}

\begin{theorem}[Proved: complete image and uniqueness, A.2(b)]
\label{C:A2b-image}
For every stable fixed pointed conformal surface and every nonempty active set $I$, the map
\[
 F_X:(0,\infty)^I\longrightarrow\mathcal K_{X,I}
\]
is a real-analytic diffeomorphism. Thus a vector of positive radii $(r_i)_{i\in I}$ is realized by a unique admissible canonical cap system, with all other ends cuspidal, if and only if
\[
 \prod_{i\in I}r_i^{a_i}
 <\prod_{a_i>0}R_i(a)^{a_i}
 \quad\text{for every }a\geq0,\quad\sum_{i\in I}a_i=1.
\]
Here $R_i(a)$ are computed from the unique Jenkins--Strebel system of circumferences $2\pi\sqrt{a_i}$. The zero radius vector corresponds to the all-cuspidal system. These statements give the whole image on all faces, and injectivity on their union.
\end{theorem}
\begin{proof}
The collar estimate puts $F_X(x)$ in $\mathcal K_{X,I}$. The energy theorem and Corollary~\ref{C:A2b-injectivity} make it injective and a local real-analytic diffeomorphism, so its image is open. We prove that the image is closed relative to $\mathcal K_{X,I}$.

Suppose $F_X(x^{(m)})\to y\in\mathcal K_{X,I}$. Compactness of the weight simplex and continuity of $H_X$ give
\[
 \delta:=\min_{a\geq0,\ \sum a_i=1}\bigl(H_X(a)-a\cdot y\bigr)>0.
\]
For all large $m$, the same gap with $y$ replaced by $F_X(x^{(m)})$ is at least $\delta/2$, uniformly in $a$. Homogeneity and Proposition~\ref{C:A2b-gap} then imply
\[
 (\delta/2)\sum_i x_i^{(m)}
 \leq H_X(x^{(m)})-x^{(m)}\cdot F_X(x^{(m)})
 \leq2g-2+n.
\]
The parameters are bounded above. No active coordinate can tend to zero, because Proposition~\ref{C:A2b-collar} would force its logarithmic radius to $-\infty$, contradicting convergence to $y$. Therefore a subsequence converges to $x$ in the positive orthant, and continuity gives $F_X(x)=y$.

The nonempty image is now both open and closed in the connected convex set $\mathcal K_{X,I}$, so it is all of it. Injectivity glues the local analytic inverses into a global inverse. Exponentiating the defining inequalities gives exactly the displayed intrinsic radius test. Different faces are distinguished by their zero entries, and the all-cuspidal face is a point. No unproved convexity assertion about arbitrary disc systems was used: the image convexity in logarithmic coordinates follows from this argument.
\end{proof}

\subsection{The full large-length limit, A.2(c)}

\begin{lemma}[Proved: compactness of normalized embedded discs]
\label{C:A2b-normal-families}
For a fixed index $i$, let $f_m:\mathbb D\to X$ be injective holomorphic maps with $f_m(0)=p_i$ whose images avoid every $p_j$, $j\ne i$. Normalize their derivatives to be positive in the chosen coordinate norm. Their conformal radii have a finite upper bound. Every sequence has a subsequence converging locally uniformly to a holomorphic map; if its radii are bounded away from zero, the limit is an injective holomorphic disc avoiding the other marked points. Nonconstant limits of mutually disjoint discs are mutually disjoint. If the radii converge to a positive value, convergence of the normalized maps to a nonconstant limit is equivalent to pointed conformal kernel convergence of their image domains.
\end{lemma}
\begin{proof}
We use Koebe's quarter and growth estimates in the form proved in Siu (2005), Appendix, pp.~10--13~\cite{C_Siu2005}. We spell out normality in all genera. If $g=0$, stability gives at least two marked points other than $p_i$. Take a fixed global coordinate sending $p_i$ to $0$ and two of the other points to $1,\infty$. The maps become univalent plane maps omitting $1$ and fixing $0$. Koebe's quarter theorem gives $|f_m'(0)|\leq4$. Its distortion bound
\[
 |f_m(w)|\leq |f_m'(0)|\frac{|w|}{(1-|w|)^2}
\]
gives local uniform boundedness. The derivative in the original coordinate differs by one fixed nonzero factor.

If $g=1$, choose a flat uniformization $\pi:\mathbb C\to X$ and normalize the lifts $\tilde f_m$ by $\tilde f_m(0)=0$. They are univalent. Since $f_m$ is an embedding and $\pi$ is locally isometric for the flat metric,
\[
 \int_{\mathbb D}|\tilde f_m'|^2\,dA
 =\operatorname{Area}_{\mathrm{flat}}(f_m(\mathbb D))
 \leq\operatorname{Area}_{\mathrm{flat}}(X).
\]
The mean-value inequality bounds $|\tilde f_m'(0)|$, and the same univalent distortion theorem gives local boundedness. If $g\geq2$, lift instead to the universal covering disc, with a fixed lift of $p_i$ at the origin; the lifted maps are bounded by one and form a normal family. In each case the normalized derivatives identify the limiting radius, and a positive limiting radius excludes a constant limit.

For completeness, the univalent-limit and disjointness assertions do not require any boundary regularity. If a nonconstant limit identified two distinct parameter points, take disjoint small parameter discs around them and a coordinate near their common image. Rouch\'e's theorem would force the nearby injective maps to take the same value on both discs, a contradiction. The same local argument excludes a critical point, an omitted marked point in the limiting image, or an intersection of two nonconstant limiting discs.

Let $D=f(\mathbb D)$ for a nonconstant normalized limit. Rouch\'e's theorem on small parameter discs shows that every compact subset of $D$ is eventually contained in $D_m=f_m(\mathbb D)$. Conversely, suppose $\Omega$ is a connected pointed open domain every compact subset of which is eventually contained in $D_m$. The inverse maps $f_m^{-1}$ are bounded holomorphic functions on those compact subsets. A diagonal normal-family argument gives $h:\Omega\to\overline{\mathbb D}$. Its derivative at $p_i$ is the positive reciprocal of the limiting radius, so it is nonconstant and $h(\Omega)\subset\mathbb D$. Passing to the limit in $f_m\circ f_m^{-1}=\operatorname{id}$ yields $f\circ h=\operatorname{id}_\Omega$. Hence $\Omega\subset D$, which proves the maximality required for kernel convergence. The reverse implication follows by the same normal-family argument and uniqueness of the normalized Riemann map.
\end{proof}

\begin{theorem}[Proved from the support-gap estimate: the complete projective limit]
\label{C:A2b-strebel-limit}
Let $x^{(m)}=(\lambda_i^{(m)})^2$ be a sequence of squared length vectors with
\[
 s_m=\sum_i x_i^{(m)}\longrightarrow\infty,
 \qquad a^{(m)}=x^{(m)}/s_m\longrightarrow a,
 \qquad a_i\geq0,\quad\sum_i a_i=1.
\]
The canonical cap at every index with $a_i>0$ converges, in the pointed conformal-kernel sense, to its unique Jenkins--Strebel disc with circumference proportional to $\sqrt{a_i}$. Every remaining cap has conformal radius tending to zero and has degenerate pointed kernel. In particular, if one boundary length tends to infinity while the others remain bounded, its cap tends to the unique one-weight Strebel slit domain. This is a domain limit on the fixed conformal surface; it does not assert Hausdorff convergence of arbitrary boundary traces.
\end{theorem}
\begin{proof}
Write $y_i^{(m)}=\log r_i^{(m)}$ for positive-length caps. Pass to a subsequence on which the set $J$ of positive-length indices is fixed; there are only finitely many such sets. Omit the indices outside $J$ throughout. Every index with $a_i>0$ belongs to $J$. Proving the asserted limit for every subsequence with fixed zero pattern proves it for the original sequence. Proposition~\ref{C:A2b-gap} gives
\[
 0\leq H_X(x^{(m)})-\sum_i x_i^{(m)}y_i^{(m)}\leq C,
 \qquad C=2g-2+n.
\]
Homogeneity gives
\begin{equation}
 \sum_i a_i^{(m)}y_i^{(m)}\geq H_X(a^{(m)})-C/s_m.
 \label{C:eq-radius-support-limit}
\end{equation}
Every coordinate has the fixed upper bound
$y_i^{(m)}\leq M_i:=H_X(e_i)<\infty$. A single fixed system of small disjoint discs also bounds $H_X$ below uniformly on the weight simplex. If $a_i>0$, then $a_i^{(m)}$ is bounded away from zero; isolating its term in~\eqref{C:eq-radius-support-limit} and using the other upper bounds gives a uniform lower bound on $y_i^{(m)}$.

Put $I=\{i:a_i>0\}$. Lemma~\ref{C:A2b-normal-families} therefore gives, along a subsequence, mutually disjoint nonconstant limiting discs $D_i$ for all $i\in I$, and convergent logarithmic radii $y_i$. Dropping the inactive terms using their upper bounds in~\eqref{C:eq-radius-support-limit}, we obtain
\[
 \sum_{i\in I}a_i y_i
 \geq\liminf_m H_X(a^{(m)})\geq H_X(a).
\]
The last inequality requires only lower semicontinuity: $H_X$ is a supremum of continuous linear functions, so it is lower semicontinuous. The weighted extremal theorem applied to the limiting disjoint active discs gives the reverse inequality. Equality follows, and its uniqueness clause identifies every $D_i$ as the designated Jenkins--Strebel disc.

These active Strebel discs cover $X$ except for their critical graph and the inactive marked points, a set with empty interior. If some inactive cap had radius bounded below along a further subsequence, Lemma~\ref{C:A2b-normal-families} would give an additional nonconstant open disc disjoint from all the active limiting discs. This is impossible. Hence every inactive radius tends to zero. No fixed open neighborhood of its center can lie eventually in such cap domains, by monotonicity of conformal radius, so their pointed kernels are degenerate.

Finally every convergent subsequence of normalized active cap maps has the same limit by uniqueness of the weighted extremal system and the derivative normalization. Normal-family compactness therefore makes the whole sequence converge. Lemma~\ref{C:A2b-normal-families} converts this to the asserted pointed kernel convergence. When only the $i$th boundary length diverges and the others stay bounded, $a=e_i$, giving precisely the one-weight slit domain.
\end{proof}

\begin{proposition}[Proved: explicit pants limits and failure of the box model]
\label{C:A2b-pants-extremes}
On $X=\mathbb P^1$ with marks $0,1,\infty$, use coordinates $z,z-1,1/z$ respectively. The one-weight extremal disc at $0$ is $\mathbb C\setminus[1,\infty)$, with radius $4$. For equal weights at $0,1$ and zero weight at $\infty$, the extremal discs are the half-planes $\operatorname{Re}z<1/2$ and $\operatorname{Re}z>1/2$, both of radius $1$. Thus every finite canonical cap system satisfies $r_0r_1<1$, although the individual supremum of each of these radii is $4$. Its image is therefore not a box. The limiting disc at $0$ is the slit domain when $\lambda_0\to\infty$ with the others bounded; it is the indicated half-plane when $\lambda_0=\lambda_1\to\infty$ and $\lambda_\infty$ stays bounded.
\end{proposition}
\begin{proof}
For one active weight take
\[
 Q=\frac{dz^2}{z^2(z-1)},\qquad
 z=\frac{4u}{(1+u)^2}.
\]
Direct substitution gives $Q=-du^2/u^2$ and the map has derivative $4$ at $0$. It maps the unit disc univalently onto $\mathbb C\setminus[1,\infty)$: its inverse is $(1-\sqrt{1-z})/(1+\sqrt{1-z})$, where the square root has positive real part off the slit, and the boundary formula $4e^{it}/(1+e^{it})^2=1/\cos^2(t/2)$ identifies its slit boundary. This supplies the one-weight characteristic disc and its radius.

For the two equal active weights, take
\[
 Q=-\frac{dz^2}{z^2(z-1)^2},\qquad u=\frac{z}{1-z}.
\]
Then $Q=-du^2/u^2$. The regions $|u|<1$ and $|u|>1$ correspond exactly to $\operatorname{Re}z<1/2$ and $\operatorname{Re}z>1/2$. Their normalized disc maps are $u\mapsto u/(1+u)$ and $v\mapsto1/(1+v)$, with derivative moduli $1$. The marked point at infinity lies on their separating trajectory. The support theorem gives $H_X(1,1,0)=0$, so the strict image inequality is $r_0r_1<1$. The one-weight computation gives $H_X(e_0)=\log4$; symmetry gives the same individual supremum at $1$. Shrinking either slit disc and inserting small discs at the other two marks shows that every radius below $4$ occurs as the corresponding coordinate of some interior admissible radius vector. Thus the simultaneous restriction cannot be a product of the individual intervals. The two rational pullback identities and radius constants are independently checked in the attached exact symbolic program. The two limits follow from Theorem~\ref{C:A2b-strebel-limit}, whose limiting squared-length weights are respectively $(1,0,0)$ and $(1/2,1/2,0)$.
\end{proof}

\paragraph{Original two-domain source.}
The same non-strict radius product bound for arbitrary disjoint simply
connected planar domains is Lavrentiev's original Theorem~3
\citep[\S2, pp.~169--173]{Lavrentiev1934}, now opened in full.
His maps $F_i:D_i\to\D$, $F_i(z_i)=0$, satisfy
$|F_0'(z_0)|\,|F_1'(z_1)|\ge |z_1-z_0|^{-2}$.
Since $r(D_i,z_i)=|F_i'(z_i)|^{-1}$, this gives
$r(D_0,z_0)r(D_1,z_1)\le|z_1-z_0|^2$.
For the canonical finite caps, our simultaneous collar enlargement
proves the strict inequality. Thus the originally remembered
Lavrentiev input is verified directly and its equality case is
distinguished from the interior admissible image.


\begin{proposition}[Cusp norm and decorations]\label{A:cuspnorm}
\status{Proved; the norm is the known Wolpert norm.}
For a marked compact Riemann surface $(X,p_1,\ldots,p_n)$ with
$2g-2+n>0$, the complete curvature $-1$ metric on
$X\setminus\{p_i\}$ defines a canonical Hermitian norm on every
$T_{p_i}X$. A normalized cusp coordinate $q$ has
\[
 ds^2=\frac{|dq|^2}{|q|^2(\log|q|)^2},\qquad
 \|\partial_q\|_{\mathrm{cusp}}=1.
\]
In its embedded range the horocycle $|q|=\rho$ has length
$h=-2\pi/\log\rho$. The unrestricted horocycle-decoration torsor is
parametrized by $h>0$, with no requirement that all projected
horocycles be disjoint. Via $h_i=\ell_i$, Theorem~\ref{A:grafting}
identifies the positive-length capping bundle with Penner's
decoration bundle. This identification does not send actual cap
shadows to cusp discs.
\end{proposition}
\begin{proof}
Uniformize a cusp by a horodisc in $\mathbb H$ with stabilizer
$w\mapsto w+1$. The coordinate $q=e^{2\pi i w}$ gives the displayed
metric by direct substitution. Any other such normalized
uniformization differs by $w\mapsto w+b$, $b\in\R$, and hence by
$q\mapsto e^{2\pi i b}q$. Thus declaring $\partial_q$ to have norm
one is independent of the choice. This is the norm of
\citet[Definition~1 and Lemma~2]{Wolpert2007}. Integrating the metric on
$|q|=\rho$ gives the formula for $h$. More generally the parabolic
cover defines a horocycle of every positive length, whether its
projection is embedded or not, yielding one positive parameter per
puncture. Penner's decoration construction \citep[\S1]{Penner1987} uses exactly these
parameters. The product descriptions therefore give the specified
bundle isomorphism and the unit-length section. The formula
$\rho=e^{-2\pi/h}$ shows why the unit tangent disc should not be
confused with a finite-length horocycle: $\rho=1$ is not an embedded
cusp neighbourhood supplied by this formula.
\end{proof}

\begin{proposition}[An exact answer to the collar metric comparison]\label{A:collar}
\status{Proved.}
Let $(X,h)$ be a compact hyperbolic surface with geodesic boundary.
Remove smooth open boundary collars so that the remaining surface
$X'$ is diffeomorphic to $X$. Let $h'$ be its canonical hyperbolic
metric with geodesic boundary. There is a unique smooth real
function $u$ on $X'$ such that $h'=e^{2u}h$, and it is characterized
by
\[
 \Delta_hu=e^{2u}-1\quad\text{on }X',\qquad
 \partial_{n}u=-k_h\quad\text{on }\partial X'.
\]
Here $\Delta_h=\operatorname{div}_h\nabla$, $n$ is the outward unit
normal, and $k_h$ is geodesic curvature measured toward the interior
of $X'$. For collars cut at constant inward distance $a_i>0$ below
their embedding widths, $k_h=-\tanh a_i$ on the corresponding new
boundary and $u>0$ throughout $X'$. Thus $h'>h$ pointwise there.
\end{proposition}
\begin{proof}
Both metrics induce the given conformal structure, so their ratio
is $e^{2u}$. The curvature transformation
$K_{e^{2u}h}=e^{-2u}(K_h-\Delta_hu)$ gives the interior equation.
The boundary transformation
$k_{e^{2u}h}=e^{-u}(k_h+\partial_nu)$ gives the Neumann equation.
Conversely, these equations imply curvature $-1$ and geodesic
boundary, and the uniqueness in bordered uniformization identifies
the resulting metric with $h'$.

For a direct uniqueness check, subtract two solutions $u,v$,
multiply by $u-v$, and integrate by parts. Their normal derivatives
agree, so
\[
 -\int_{X'}|\nabla(u-v)|^2\,dA_h
 =\int_{X'}(u-v)(e^{2u}-e^{2v})\,dA_h\geq0.
\]
Hence $u=v$. In Fermi coordinates $h=da^2+\cosh^2(a)\,ds^2$;
the curve $a=a_i$, with the remaining surface on its increasing-$a$
side, has inward geodesic curvature $-\tanh a_i$.
Thus $\partial_nu=\tanh a_i>0$. Write
$\Delta_hu-c(x)u=0$, where
$c(x)=(e^{2u(x)}-1)/u(x)>0$ with value $2$ at $u=0$.
The maximum principle excludes a negative interior minimum;
the boundary point lemma excludes a nonpositive boundary minimum
because its outward derivative would have to be negative.
The solution is not identically zero since the prescribed normal
derivatives are positive. The strong maximum principle gives
$u>0$. For arbitrary non-equidistant collars the displayed boundary
value problem remains the exact comparison; no unproved monotonicity
is asserted.
\end{proof}

% Independent completion of the unrestricted collar question, by /root/algebra_c.
% AI-produced by GPT-6 Astra (OpenAI), ChatGPT Pro mode, 19 November 2026.
\subsection{Arbitrary collars: there is no pointwise metric monotonicity}
\label{C:A2e-section}

The official scan~67/021 asks for the relation between the canonical metrics after removing an open ``collar''; it does not require that the new boundary be an equidistant curve. Proposition~\ref{A:collar} gives the exact nonlinear comparison and the strict increase for equidistant collars. The following result settles the unrestricted smooth, and even real-analytic, case.

\begin{theorem}[Proved: arbitrarily small analytic collars give incomparable metrics]
\label{C:A2e-nonmonotonicity}
Let $(X,h)$ be any compact connected hyperbolic surface with nonempty geodesic boundary and negative Euler characteristic. There are arbitrarily thin open collars of its boundary, with real-analytic inner boundaries, such that the remaining surface $X_t$ is diffeomorphic to $X$ but its canonical curvature $-1$ metric $h_t'$ with geodesic boundary is neither pointwise greater than nor pointwise less than $h|_{X_t}$. More precisely, their conformal ratio satisfies
\[
 h_t'=e^{2u_t}h|_{X_t},
 \qquad \min_{X_t}u_t<0<\max_{X_t}u_t.
\]
This includes every compact hyperbolic pair of pants and every compact one-holed torus. The new boundaries can converge to the original ones in $C^m$ for every fixed $m$.
\end{theorem}

\begin{proof}
Fix one boundary component $B$, of length $L$, and a disjoint Fermi collar $0\le r<R$ around it, with arclength coordinate $s\in\mathbb R/L\mathbb Z$. Here
\[
 h=dr^2+\cosh^2r\,ds^2,
 \qquad n|_B=-\partial_r.
\]
Choose $0<\varepsilon<1$, and, for a positive integer $N$ to be chosen, put $k=2\pi N/L$ and
\[
 f(s)=1+\varepsilon\cos(ks)\quad\text{on }B,
 \qquad f=1\quad\text{on every other boundary component}.
\]
Remove the relatively open collars $0\le r<t f(s)$, using Fermi coordinates at every boundary component. Since $f\ge1-\varepsilon>0$, these are collars of strictly positive width. For each fixed $N$ and all sufficiently small $t>0$, their closures are disjoint, their inner boundaries are analytic, and their complement $X_t$ is diffeomorphic to $X$.

We first derive the linear response rigorously, including the dependence on the moving domain. Choose smooth diffeomorphisms $T_t:X\to X_t$ which take $(0,s)$ to $(tf(s),s)$ and are the identity outside the fixed collars; for example, near a boundary take $r\mapsto r+t\chi(r)f(s)$ with $\chi=1$ near $0$ and $\chi=0$ farther inside. Put $\widehat h_t=T_t^*h$ and $U_t=u_t\circ T_t$. Then $\widehat h_t\to h$ smoothly, with difference $O(t)$ in every fixed differentiability norm. Proposition~\ref{A:collar} becomes
\begin{equation}
 \Delta_{\widehat h_t}U_t=e^{2U_t}-1,
 \qquad \partial_{n_t}U_t=-\kappa_t,
 \label{C:eq-collar-pulled}
\end{equation}
where $\kappa_t$ is the $h$-geodesic curvature of $r=tf(s)$ toward the remaining surface, pulled back to the original boundary. Direct differentiation in the displayed Fermi metric gives
\begin{equation}
 \kappa_t=t(f''-f)+O(t^2)
 \label{C:eq-collar-curvature}
\end{equation}
in every fixed differentiability norm. Indeed the radial component of the acceleration of $s\mapsto(tf(s),s)$ is $t f''-\cosh(tf)\sinh(tf)=t(f''-f)+O(t^2)$, while its unit tangent and inward unit normal have the limiting values $\partial_s$ and $\partial_r$; the tangential acceleration contributes only order $t^2$ to the curvature.

For clarity, the linear boundary problem used below is
\begin{equation}
 (\Delta_h-2)w=0,\qquad \partial_nw=b.
 \label{C:eq-collar-linear}
\end{equation}
It has a unique smooth solution for every smooth boundary function $b$. Its weak formulation is
\[
 \mathcal E(w,v):=\int_X
       (\langle\nabla w,\nabla v\rangle_h+2wv)\,dA_h
   =\int_{\partial X}bv\,ds_h.
\]
The trace theorem makes the right side bounded on $H^1(X)$, and $\mathcal E$ is coercive there, so the Hilbert-space representation theorem gives existence and uniqueness. Elliptic boundary regularity gives smoothness. The same linear problem, including its existence and estimate, appears in Rupflin (2021), Lemma~2.9(ii), equation~(2.22), pp.~3126--3127, specialized to geodesic boundary~\cite{C_Rupflin2021}. Unlike a pure Neumann Laplace equation, the term $-2w$ removes the constant-function kernel, so no compatibility condition is needed.

We show that the solution of~\eqref{C:eq-collar-pulled} obeys
\begin{equation}
 U_t/t\longrightarrow w_f\quad\text{in }H^1(X),
 \qquad \partial_nw_f=f-f''.
 \label{C:eq-collar-response}
\end{equation}
This argument does not presume a shape derivative theorem. Let $w_0$ solve~\eqref{C:eq-collar-linear} with $b=1$ on the entire boundary. The maximum principle, or the negative-part test in its weak equation, gives $w_0>0$. Set $P=w_0+1$. For a sufficiently large fixed constant $A$ and sufficiently small $t>0$, the functions $\pm AtP$ are upper and lower barriers for~\eqref{C:eq-collar-pulled}. To verify this explicitly, use
\[
 (\Delta_{\widehat h_t}-2)P=-2+O(t),
 \qquad \partial_{n_t}P=1+O(t),
 \qquad \kappa_t=O(t).
\]
Thus $\Delta_{\widehat h_t}(AtP)-(e^{2AtP}-1)<0$, its outward derivative is at least $-\kappa_t$, and the reversed inequalities hold for $-AtP$. Comparison follows by multiplying the equation for the positive part of the difference by that positive part and integrating; the interior zeroth-order coefficient is strictly positive and the boundary term has the opposite sign. Consequently $\|U_t\|_\infty\le Ct$.

Let $\mathcal E_t$ be the preceding energy form with metric $\widehat h_t$. The weak equation is
\[
 \mathcal E_t(U_t,v)
 =-\int_{\partial X}\kappa_t v\,ds_{\widehat h_t}
  -\int_X(e^{2U_t}-1-2U_t)v\,dA_{\widehat h_t}.
\]
Since $\|U_t\|_\infty=O(t)$, the second integrand is uniformly $O(t^2)|v|$. Uniform coercivity, the trace theorem, and the choice $v=U_t$ give $\|U_t\|_{H^1}=O(t)$. Divide the equation by $t$, subtract the weak equation with datum $f-f''$, and use~\eqref{C:eq-collar-curvature}. The differences of the metric coefficients and measures are $O(t)$, and the exponential remainder divided by $t$ is $O(t)$. We obtain
\[
 |\mathcal E(U_t/t-w_f,v)|\le Ct\|v\|_{H^1}.
\]
Taking $v=U_t/t-w_f$ proves~\eqref{C:eq-collar-response}, in fact with error $O(t)$. In particular, its boundary traces converge in $L^2$.

We now choose $N$ so that this first-order response is negative somewhere. Let $w_k$ solve~\eqref{C:eq-collar-linear} with datum $(1+k^2)\cos(ks)$ on $B$ and zero on all other boundaries. Then
\[
 w_f=w_0+\varepsilon w_k.
\]
Let $\phi_B$ solve the same linear problem with datum $1$ on $B$ and zero elsewhere. Reciprocity of $\mathcal E$ gives the exact identities
\begin{align}
 (1+k^2)\int_Bw_k\cos(ks)\,ds
   &=\mathcal E(w_k,w_k),\label{C:eq-collar-fourier-energy}\\
 \int_Bw_k\,ds
   &=(1+k^2)\int_B\phi_B\cos(ks)\,ds.
   \label{C:eq-collar-mean}
\end{align}
Set $K_3=\int_B|\partial_s^3\phi_B|\,ds<\infty$. Three integrations by parts on the closed boundary circle show that, for $k\ge1$,
\begin{equation}
 \left|\int_Bw_k\,ds\right|
 \le (1+k^2)K_3/k^3\le2K_3/k.
 \label{C:eq-collar-mean-bound}
\end{equation}

For the opposite estimate on its oscillating coefficient, use the explicit $H^1$ test function
\[
 v_k(r,s)=(1-r/R)e^{-kr}\cos(ks)\quad(0\le r\le R),
 \qquad v_k=0\quad\text{outside this collar}.
\]
It is continuous at $r=R$, so its zero extension lies in $H^1(X)$. Put $\chi(r)=1-r/R$. Direct integration of the angular variables gives
\[
\begin{aligned}
 \mathcal E(v_k,v_k)
 &=\frac L2\int_0^R e^{-2kr}
       \bigl[\cosh r\,(\chi'-k\chi)^2\\
 &\hspace{6em}+(k^2\operatorname{sech}r+2\cosh r)\chi^2\bigr]\,dr.
\end{aligned}
\]
For $k\ge\max(1,R^{-1})$, the bounds $(a+b)^2\le2a^2+2b^2$, $0\le1-r/R\le1$, and $\int_0^R e^{-2kr}dr\le(2k)^{-1}$ yield
\[
 \mathcal E(v_k,v_k)\le C_R k,
 \qquad C_R=\frac{L(6\cosh R+1)}4.
\]
But $v_k|_B=\cos(ks)$, and hence
\[
 \mathcal E(w_k,v_k)=(1+k^2)L/2.
\]
Cauchy--Schwarz for the positive inner product $\mathcal E$, followed by~\eqref{C:eq-collar-fourier-energy}, proves
\begin{equation}
 \int_Bw_k\cos(ks)\,ds
 \ge\frac{(1+k^2)L^2}{4C_Rk}
 \ge c_Rk,
 \qquad c_R=\frac{L}{6\cosh R+1}>0.
 \label{C:eq-collar-amplitude}
\end{equation}
No asymptotic expansion of a boundary pseudodifferential operator is used.

Let $M_0=\int_Bw_0\,ds>0$. Since $w_0>0$, its cosine coefficient has absolute value at most $M_0$. Equations~\eqref{C:eq-collar-mean-bound} and~\eqref{C:eq-collar-amplitude} imply
\[
 \int_B w_f(1-\cos(ks))\,ds
 \le2M_0+2\varepsilon K_3/k-\varepsilon c_Rk.
\]
Choose the integer $N$ so that $k=2\pi N/L$ satisfies
\[
 k>\max\left(1,R^{-1},\frac{4M_0}{\varepsilon c_R},
                        \sqrt{\frac{4K_3}{c_R}}\right).
\]
The displayed integral is then strictly negative. Its weight $1-\cos(ks)$ is nonnegative, so $w_f$ is negative at some point of $B$.

By~\eqref{C:eq-collar-response} and trace convergence,
\[
 \int_B U_t(1-\cos(ks))\,ds<0
\]
for all sufficiently small positive $t$. Therefore $U_t<0$ at some boundary point. Smoothness gives $u_t<0$ also on a nonempty relatively open region of $X_t$, including interior points, proving $h_t'<h$ there.

There must also be a region with the opposite inequality. Gauss--Bonnet and the fact that both canonical metrics have geodesic boundary give
\[
 \operatorname{Area}_{h_t'}(X_t)
  =-2\pi\chi(X_t)=-2\pi\chi(X)
  =\operatorname{Area}_h(X)
  >\operatorname{Area}_h(X_t).
\]
Hence $\int_{X_t}(e^{2u_t}-1)\,dA_h>0$, so $u_t>0$ somewhere. This proves incomparability. Finally $N$ is fixed before $t\downarrow0$; all derivatives of $tf(s)$ tend to zero. Thus the collars are arbitrarily thin and their analytic boundaries converge in every $C^m$ norm as asserted.
\end{proof}

\paragraph{Exact scope of the answer to 67/021.}
The nonlinear Neumann equation in Proposition~\ref{A:collar} characterizes the comparison for any smooth collar. Equidistant inward cuts strictly increase the metric. General analytic inward cuts can make the metrics incomparable, by Theorem~\ref{C:A2e-nonmonotonicity}. The manuscript itself asks a question and does not assert a false monotonicity theorem. Our counterexample refutes that potential generalization, while retaining the stated hypotheses of the question.


\section{A.3. Correct welding statements}

\begin{theorem}[Welding with an explicit regularity class]\label{A:welding}
\status{Known existence theorem, with the category and uniqueness specified.}
Let compact oriented bordered Riemann surfaces have finitely many
paired boundary circles, and glue each pair by an orientation-reversing
quasisymmetric homeomorphism, measured in analytic boundary
coordinates. The topological quotient carries a conformal structure
which restricts to the given structures on every interior, and its
seams are quasicircles. It is unique under the identity on the
pieces. More generally, for any given welding homeomorphisms,
two realizations are conformally identical if one realization's
seams are conformally removable. Rectifiable seams satisfy this
uniqueness condition. Quasisymmetry is sufficient for existence
and characterizes quasicircle realizations; it is not necessary
for existence of arbitrary Jordan-curve weldings.
\end{theorem}
\begin{proof}
The quasisymmetric extension theorem extends each gluing map
quasiconformally into a collar. Choose disjoint collars, build the
corresponding quasiconformal quotient atlas, and solve its bounded
Beltrami coefficient. This yields a Riemann surface with the
prescribed conformal interiors and quasiconformal images of circles
as seams. The atlas construction is invariant under changing
analytic collar coordinates, since those changes are conformal.

For uniqueness, the identity on corresponding pieces is a global
homeomorphism between two realizations, conformal off the seams.
Conformal removability says exactly that this homeomorphism is
conformal everywhere. Quasicircles are conformally removable; a
rectifiable curve is removable for continuous holomorphic
functions, so it is also removable for such homeomorphisms.
The converse for the quasicircle class follows by extending its
two conformal maps quasiconformally across the circle and taking
their boundary composition. No converse for arbitrary Jordan
seams is used.
The planar construction is given explicitly in
\citet[Section 2]{Hamilton2002}; \citet[Section 1]{Younsi2018}
states the removability implication and its precise scope.
\end{proof}

\begin{proposition}[Slits and the corrected Lemma 4B and Scholie]\label{A:slit}
\status{Proved from the welding theorem and prime-end uniformization.}
Let $V$ be a Jordan disc and let $I$ be a Jordan arc with one
endpoint on $\partial V$ and its remaining points in $V$. Cutting
$V$ along $I$ gives a simply connected surface with a prime-end
boundary circle. Its Riemann map has the boundary identifications
of 67/071--073, with the slit tip and its two access points at
the attachment prescribed as the three normalization points.
No rectifiability is needed for this cutting assertion.

Conversely, let the two designated boundary banks be identified
by $\varphi$, and form their bank-exchanging involution $H$:
$H=\varphi$ on one bank and $H=\varphi^{-1}$ on the other.
Require $H$ to be quasisymmetric in an analytic boundary
coordinate on the \emph{union} of the banks, including their
common endpoint. This condition has a slit realization which
is the image of a straight radial slit under a quasiconformal
homeomorphism of a disc. After putting its tip at $0$ and its
attachment at $1$ in the unit disc, the normalized realization
is unique. Conversely every such quasiconformal slit pair has
this welding condition. The correspondence is therefore with
this specified class of slit pairs; quasisymmetry on a single
bank alone is insufficient at the fold.
For a prescribed rectifiable slit the normalized welding is
unique, but ``rectifiable graph'' does not supply an existence
criterion for an arbitrary interval homeomorphism.
\end{proposition}
\begin{proof}
The cut-open Jordan disc is topologically a closed disc. Its
boundary is locally connected, including the two banks and the
slit tip. Prime-end uniformization and Caratheodory's theorem
give a continuous boundary parametrization with precisely the
specified identifications. The three prime ends determine its
disc automorphism uniquely because $\operatorname{Aut}(\D)$
acts simply transitively on positively ordered boundary triples.

For the converse put the common bank endpoint at $0$ and the
other bank endpoints at $-1,1$ in a boundary interval. Thus
$H$ is decreasing, $H^2=\mathrm{id}$, $H(0)=0$, and
$H(\pm1)=\mp1$. Define
\[
 k(x)=\begin{cases}x&x\geq0,\\-H(x)&x\leq0.\end{cases}
\]
Then $kHk^{-1}(x)=-x$. The restriction of $k$ to each
half-interval is quasisymmetric. Across zero, quasisymmetry
of $H$ gives $|H(-t)|/t$ bounded above and below: write $a(t)=|H(t)|$, so $H(-t)=a^{-1}(t)$.
The equal-adjacent-interval condition gives
$K^{-1}\le a(t)/a^{-1}(t)\le K$. If $a(t)\ge t$, then
$a^{-1}(t)\le t$ and $a(t)\le Kt$; the reverse case gives
$a(t)\ge t/K$. Thus both $a(t)$ and $a^{-1}(t)$ are
comparable to $t$.
Equivalently one may use the two explicit conditions of
\citet[Lemma 6]{Lind2005}: one-bank quasisymmetry and
comparability of opposite distances from the fold. Splitting
an interval crossing zero into at most three adjacent
intervals then proves the same condition for $k$. Hence $k$
extends to a quasisymmetric boundary-circle homeomorphism:
reflect it near $-1,1$ on the remaining boundary arc and
interpolate on the remaining compact interval. The reflection
makes the endpoint estimates match, and finitely many compact
interval estimates give a uniform constant.

Extend this boundary map quasiconformally to the disc.
Composing with the standard quotient which folds two equal
banks to a radial slit gives a quasiconformal map on the
interior of the cut disc, with the prescribed identifications
on its boundary. Transfer its complex structure to the
model disc minus the straight slit. Its Beltrami coefficient
is bounded strictly below one; extend it arbitrarily over
the slit, a set of area zero. The measurable Riemann mapping
theorem straightens this structure. The resulting quotient
map is holomorphic on the original open disc, and the seam
is the quasiconformal image of the straight slit. Uniformize
the quotient disc, taking the tip to $0$ and attachment to
$1$. Removability of the seam gives uniqueness, because a
disc automorphism fixing $0$ and $1$ is the identity.
For the converse, open an ambient quasiconformal map carrying
the straight slit pair to the prescribed slit pair. Its
boundary map on the cut-open discs is quasisymmetric, so it
conjugates reflection to a quasisymmetric involution $H$.
The same uniqueness argument applies to any rectifiable slit.

Finally, every increasing interval homeomorphism has bounded
variation, and its graph has length at most the sum of its
horizontal and vertical variations. Hence all such graphs
are rectifiable. That graph condition cannot distinguish
the analytic welding classes above. On the original scan
67/073 the first normalization is $f(\omega)=0$: $\omega$
is the slit-tip prime end. This resolves the transcription's
uncertain alternative $\omega\mid0$ at its line 745. The alternative interior
reading would be impossible, since the interior maps onto
the slit complement and cannot contain the tip.
\end{proof}

\begin{example}[Why the fold condition cannot be suppressed]\label{A:fold}
\status{Proved counterexample to the weaker quasisymmetric formulation.}
The interval homeomorphism $h(t)=t^2$ on $[0,1]$ is
quasisymmetric, whereas the decreasing involution
\[
 H(x)=\begin{cases}-x^2&x\geq0,\\\sqrt{-x}&x\leq0\end{cases}
 \quad(-1\leq x\leq1)
\]
is not. For equal adjacent intervals at $0$, the ratio of
the image lengths is $t^{1/2}/t^2=t^{-3/2}$, unbounded as
$t\downarrow0$. To check the first assertion directly, for
$t>0$ and $x\geq t$ the adjacent-increment ratio for $h$ is
$(2x+t)/(2x-t)\in[1,3]$; the usual interval definition
only tests triples in $[0,1]$. Thus the datum is
quasisymmetric on one bank but cannot be the welding of a
quasiconformal slit pair. We make no assertion here that it
has no Jordan-slit realization in a larger class.
\end{example}

\begin{corollary}[The scope of L5]\label{A:L5}
\status{Proved under the stated geodesic-cut convention.}
On 67/074--082, geodesic cuts and the rotational capping
identifications are real analytic, hence quasisymmetric and
removable. No earlier assertion about arbitrary rectifiable
weldings is needed for that construction. If ``differentiable
cut'' is taken literally without toric compatibility, the
obstruction in Proposition~\ref{A:arbitrarycuts} remains;
smoothness alone does not produce the claimed nodal tensor.
\end{corollary}
\begin{proof}
An analytic circle diffeomorphism has derivative bounded
above and away from zero on compact coordinate arcs, so is
locally bi-Lipschitz and quasisymmetric. Geodesic seams
have finite length and are removable. The construction
therefore uses only the valid case of the welding theorem.
The last assertion is exactly the explicit nonrotational
gluing in Proposition~\ref{A:arbitrarycuts}.
\end{proof}

\section{A.4. The letter to Bers and its intrinsic replacement}

\begin{theorem}[The orange-slice model fails]\label{A:orange}
\status{Refuted; the method is credited to Michel Hua's project.}
The model proposed on 67/019--020, in which all admissible pants caps
are concentric subdiscs of three fixed orange slices, is false.
There is a hyperbolic pair of pants with two cusps and third
boundary length $32$ whose canonical cap cannot lie in the
proposed cell at $0$.
\end{theorem}
\begin{proof}
Normalize the cusp centres to $1,\infty$ and the cap centre to $0$.
In these coordinates the proposed orange cell at zero is
$Q_0=\{|z|\leq1,\ \operatorname{Re}z\leq1/2\}$. This is the cell
of the round metric for which $0,1,\infty$ lie equally spaced on a
great circle. Its finite boundary pieces lie on $|z|=1$ and
$\operatorname{Re}z=1/2$.

For every $z\in Q_0$,
$|z+1/2|^2=|z|^2+\operatorname{Re}z+1/4\leq7/4<(4/3)^2$.
The puncture $1$ has distance $3/2>7/5$ from $-1/2$.
Consequently, if $D_0\subset Q_0$, the Euclidean annulus
\[
 A=\{4/3<|z+1/2|<7/5\}
\]
lies in $P=\PP^1\setminus(D_0^\circ\cup\{1,\infty\})$
and its core curves are freely homotopic to the geodesic boundary
of $P$. Write $\operatorname{Ext}_P(\gamma)$ for the extremal
length of this class. Restriction to the annulus gives
\[
 \operatorname{Ext}_P(\gamma)\leq
 \operatorname{Ext}_A(\gamma)
 =\frac{2\pi}{\log(21/20)}.
\]
The curvature $-1$ metric on $P$ has area $2\pi$ by Gauss--Bonnet,
and the shortest representative of this class has length $L$.
Using that metric in the definition of extremal length gives
$\operatorname{Ext}_P(\gamma)\geq L^2/(2\pi)$.
For $L=32$ these bounds contradict each other, since
\[
 32^2\log(21/20)>4\pi^2.
\]
For a completely elementary numerical verification,
$\log(1+x)>x/(1+x)$ for $x>0$ gives
$1024\log(21/20)>1024/21>40$, while
$4\pi^2<4(22/7)^2<40$.
The required hyperbolic pants exists for the length triple
$(32,0,0)$, with zero lengths interpreted as cusps. Its canonical
capping therefore gives the asserted counterexample. This proof
uses Hua's extremal-length obstruction and verifies the particular
annulus and numerical constants independently.
\end{proof}

\begin{corollary}[The intrinsic multiplicative action]\label{A:action}
\status{Proved from Theorem~\ref{A:grafting}.}
The semigroup $[0,\infty)^I$ acts continuously on
$\mathcal T^{B,\geq0}(S)$ by
\[
 a\cdot Y=\overline\Psi^{-1}
 \bigl(\operatorname{cap}(Y),(a_i\ell_i(Y))_{i\in I}\bigr).
\]
It is intrinsic and mapping-class equivariant; it is real analytic
on each stratum on which the source and target patterns of zero
lengths are fixed. The group $(0,\infty)^I$ acts freely and
transitively on each positive-length fibre. The all-zero multiplier
is the capping retraction. The quotient is $\mathcal T_{g,n}$.
The section with all lengths equal to $1$ is canonical.
\end{corollary}
\begin{proof}
In the product coordinates of Theorem~\ref{A:grafting}, all claims
are the corresponding elementary claims for componentwise
multiplication on $[0,\infty)^I$. Transporting the operation by the
specified bijection proves the action law and continuity. The map
$\Psi$ is natural under marked conformal isomorphisms and mapping
classes, so no pants decomposition enters the definition. For two
positive length vectors in the same fibre, the unique multiplier
is their componentwise ratio. Multiplication by zero fixes precisely
the zero section. Projection from the product is the quotient map;
this also proves the asserted quotient topology.
\end{proof}

